% This class has a lot of options, so please check deepmind.cls for more details.
% This is a minimal set for most needs.
\documentclass[11pt, a4paper, logo, copyright]{googledeepmind}

% Omit dates for reproducibility.
\pdfinfoomitdate 1
\pdftrailerid{redacted}

% This avoids duplicate hyperref bookmark entries when using \bibentry (e.g. via \citeas).
\makeatletter
\renewcommand\bibentry[1]{\nocite{#1}{\frenchspacing\@nameuse{BR@r@#1\@extra@b@citeb}}}
\makeatother

\usepackage{kantlipsum, lipsum}
\usepackage{dsfont}
\usepackage{gdm-colors}
\usepackage[utf8]{inputenc}   % 允许 UTF-8 输入
\usepackage{newunicodechar}   % 定义 Unicode 字符
\usepackage{amssymb}          % 提供部分符号支持
\usepackage{arydshln}
\usepackage{wasysym,marvosym}
 
% Sometimes you will get errors about pdflink ending up in diffrent position. Try this and
% comment it out again when you are done with your document.
%\hypersetup{draft}

% Set the bibliography options here.
\usepackage[authoryear, sort&compress, round]{natbib}

\let\Cross\undefined
\let\Square\undefined
\usepackage{bbding}
\usepackage[utf8]{inputenc} % allow utf-8 input
\usepackage[T1]{fontenc}    % use 8-bit T1 fonts
\usepackage{hyperref}       % hyperlinks
\usepackage{url}            % simple URL typesetting
\usepackage{booktabs}       % professional-quality tables
\usepackage{nicefrac}       % compact symbols for 1/2, etc.
\usepackage{microtype}      % microtypography
\usepackage{amsmath}
\usepackage{graphicx}
\usepackage{tablefootnote}
\usepackage{multicol}
\usepackage[nameinlink]{cleveref}
\usepackage{bbm}
\usepackage{multirow}
% \usepackage{subfig}
\usepackage{soul}
% \usepackage{floatrow}
\usepackage{float}
\usepackage{wrapfig}
\usepackage{blindtext}
\usepackage{tablefootnote}
% \usepackage{fdsymbol}. % cause an error
\usepackage{amsfonts}
\usepackage[flushleft]{threeparttable}
\usepackage{colortbl}
\usepackage{mathtools,amssymb}
\usepackage{bm}
\usepackage{makecell}
\usepackage{caption}
\usepackage{capt-of}
\usepackage{array}
\usepackage{calc}      % 用于测量高度
\usepackage{caption}   % 用于 \captionof 命令
% \usepackage{subfigure} % 用于创建子图
\usepackage{subcaption}  % 在导言区添加
\usepackage{xcolor,colortbl}
\usepackage[bottom]{footmisc}
\newcommand{\ruiming}[1]{\textcolor{red}{#1}} 
\newcommand{\yunfan}[1]{\textcolor{yellow}{#1}} 
\newcommand{\jinghao}[1]{\textcolor{olive}{#1}} 



\usepackage{tcolorbox}

\definecolor{mydeepgreen}{RGB}{0, 100, 0} 

\newcommand{\myred}[1]{\textcolor{red}{#1}}
\newcommand{\myblue}[1]{\textcolor{blue}{#1}}      % 蓝色
\newcommand{\mygreen}[1]{\textcolor{mydeepgreen}{#1}}    % 绿色
\newcommand{\myyellow}[1]{\textcolor{yellow}{#1}}  % 黄色
\newcommand{\mypurple}[1]{\textcolor{purple}{#1}}  % 紫色
% \newcommand{\myorange}[1]{\textcolor{orange}{#1}}  % 橙色

\newcommand{\argmin}{\operatornamewithlimits{argmin}}
\newcommand{\argmax}{\operatornamewithlimits{argmax}}




%%%%%%%%%%% mboratko: added these, feel free to adjust
% \DeclareMathOperator*{\argmax}{arg\,max}
% \DeclareMathOperator*{\argmin}{arg\,min}
%%%%%%%%%%


\input{commands}

% Images will be looked for in this path, removes need for explicit path when including images.
\graphicspath{{figures/}}

% Important Information about your paper.
% \title{Kelix: Closing the Understanding Gap of Discrete Tokens in Unified Multimodal Models}

\title{Kelix Technique Report: \\ Closing the Understanding Gap of Discrete Tokens in Unified Multimodal Models}

% Override the running header to show short title only
\fancyhead[L,R]{}
\fancyhead[C]{\footerfont Kelix Technique Report}

% Can leave this option out if you do not wish to add a corresponding author.
%\correspondingauthor{jinhyuklee@google.com}

% Remove these if they are not needed
% \keywords{\LaTeX, Publications process, tools}
% \paperurl{arxiv.org/abs/123}

% Use the internally issued paper ID, if there is one
% \reportnumber{001} % Leave blank if n/a

% Assign your own date to the report.
% Can comment out if not needed or leave blank if n/a.
\renewcommand{\today}{}

% Can have as many authors and as many affiliations as needed. Best to indicate joint
% first-authorship as shown below.
\author{\large OneRec Team}

% \author[ ]{Yi Luan\hspace{-0.4ex}}
% \author[ ]{Sai Meher Karthik Duddu\hspace{-0.4ex}}
% \author[ ]{Gustavo Hernandez Abrego\hspace{-0.4ex}}
% \author[ ]{Weiqiang Shi\hspace{-0.4ex}}
% \author[ ]{Nithi Gupta\hspace{-0.4ex}}
% \author[ ]{Aditya Kusupati\hspace{-0.4ex}}
% \author[ ]{Prateek Jain\hspace{-0.4ex}}
% \author[ ]{Siddhartha Reddy Jonnalagadda\hspace{-0.4ex}}
% \author[ ]{Ming-Wei Chang\hspace{-0.4ex}}
% \author[ ]{Iftekhar Naim}

% Jinhyuk Lee, Zhuyun Dai, Xiaoqi Ren, Blair Chen, Daniel Cer, Jeremy R. Cole, Kai Hui, Michael Boratko, Rajvi Kapadia, Wen Ding, Yi Luan, Sai Meher Karthik Duddu, Gustavo Hernandez Abrego, Weiqiang Shi, Nithi Gupta, Aditya Kusupati, Prateek Jain, Siddhartha Reddy Jonnalagadda, Ming-Wei Chang, Iftekhar Naim

% Affiliations *must* come after the declaration of \author[]
%\affil[*]{Equal contributions}
% \affil[ ]{Google}
% \affil[2]{Google Cloud}
% \affil[3]{Google Research}
% \affil[4]{Google Labs}

\begin{abstract}
    Autoregressive large language models (LLMs) scale well by expressing diverse tasks as sequences of discrete natural-language tokens and training with next-token prediction, which unifies comprehension and generation under self-supervision. Extending this paradigm to multimodal data requires a shared, discrete representation across modalities.
    %
    However, most vision--language models (VLMs) still rely on a hybrid interface: discrete text tokens paired with continuous Vision Transformer (ViT) features. Because supervision is largely text-driven, these models are often biased toward understanding and cannot fully leverage large-scale self-supervised learning on non-text data.
    %
    Recent work has explored discrete visual tokenization to enable fully autoregressive multimodal modeling, showing promising progress toward unified understanding and generation. Yet existing discrete vision tokens frequently lose information due to limited code capacity, resulting in noticeably weaker understanding than continuous-feature VLMs.
    %
    We present \myorange{Kelix}, a fully discrete autoregressive unified model that closes the understanding gap between discrete and continuous visual representations. At its core is a novel multi-token vision tokenization mechanism that increases the information capacity of discrete codes, enabling discrete-token models to achieve understanding performance comparable to their continuous-feature counterparts. Extensive experiments demonstrate state-of-the-art results on most multimodal understanding and image generation benchmarks among comparable-scale unified models. In particular, Kelix achieves significant improvements over prior discrete approaches on understanding tasks, especially on text-rich benchmarks, where discrete methods have been weakest, surpassing the previous best discrete model by over 23\% on OCRBench and closely approaching continuous baselines. This demonstrates that discrete representations are no longer the bottleneck of multimodal understanding, offering new insights for the future development of discrete unified models.
\end{abstract}






\begin{document}

\maketitle


\begin{figure}[h!]
    \centering
    \includegraphics[width=1\textwidth]{OneBase/figures/rader.pdf}
      \caption{Benchmark performance of Kelix. Compared with prior discrete unified models, Kelix significantly improves understanding, even matching continuous understanding-only VLMs, while achieving competitive image generation results with current SOTA.}
      \label{rader}
\end{figure}
  


  







\section{Introduction}


Autoregressive large language models (LLMs) owe much of their scaling success to a unified discrete representation: by expressing diverse tasks as sequences of tokens and training with next-token prediction, a single objective serves both comprehension and generation~\citep{radford2018improving,radford2019language,brown2020language,achiam2023gpt,agarwal2025gpt,kaplan2020scaling}. Extending this paradigm to multimodal data, where vast non-text corpora remain underexploited, promises a new scaling frontier.
%
Yet mainstream vision-language models (VLMs) still rely on a hybrid interface: discrete text tokens paired with continuous Vision Transformer (ViT) features~\citep{liu2023visual,dosovitskiy2020image}. Because supervision is largely text-driven, vision often collapses into a conditioning signal, yielding understanding-only models that cannot fully leverage self-supervised learning on non-text data~\citep{yao2024minicpm,chen2024internvl,team2025kimi}.
%
A natural solution is to align vision to the discrete token interface of LLMs via discrete image tokenizers, enabling fully autoregressive multimodal modeling that unifies understanding and generation. However, existing discrete-token approaches suffer from limited information capacity: a single discrete code represents far less information than the continuous embedding it replaces, resulting in noticeably weaker multimodal understanding than continuous-feature VLMs.

In this paper, we propose \textbf{Kelix}\footnote{The name \emph{Kelix} is a portmanteau of \emph{K}uaishou and h\emph{elix}: just as the DNA double helix encodes the full complexity of life using only four discrete nucleotide bases (A, T, C, G), Kelix encodes rich visual semantics using discrete tokens.}, a fully discrete, LLM-centric unified model that closes this understanding gap. At its core is a multi-token vision tokenization scheme that expands the coding space of discrete visual representations, substantially reducing information loss. The tokenizer is learned end-to-end with the LLM, allowing quantization to adapt to the downstream objective. We further adopt a next-block prediction objective that models grouped discrete token blocks, mitigating context-length explosion while preserving multimodal information. Extensive experiments demonstrate state-of-the-art results on major understanding and generation benchmarks among comparable-scale unified models.










\section{Design Philosophy}
\label{sec2_Design}

Before describing the architecture in detail, we outline the design principles that guided our choices. Unified multimodal modeling has not yet converged on a single dominant design. Rather than exploring the full space empirically, we take \emph{representation design} as the entry point: how should input/output token representations be defined so that modalities can be aligned and fused within a single LLM-centric framework?

We factor representation design into two largely orthogonal dimensions:
\begin{itemize}
    \item \myorange{\textbf{Representation form}}: discrete vs.\ continuous.
    \item \myorange{\textbf{Semantic level}}: high-level semantics vs.\ low-level semantics.
\end{itemize}
Our choices are guided by two principles: (1)~\textbf{strong scaling potential}---the representation should support efficient scaling via high sample efficiency, high MFU, and a unified space for cross-modal transfer; and (2)~\textbf{maximum compatibility with the LLM ecosystem}---reusing open foundation models, training frameworks, and deployment toolchains.

\textbf{Semantic level.} We prefer \emph{high-level semantics} as the unified representation. High-level tokens are more compressible and align naturally with text, benefiting sample efficiency and cross-modal transfer. To maximize MFU, we decouple the system into a modular \emph{Tokenizer--LLM--Detokenizer} pipeline: the LLM backbone performs scalable sequence modeling in the semantic space, while modality-specific encoding and decoding are delegated to the Tokenizer and Detokenizer, which interact with the LLM only through the unified representation interface.

\textbf{Representation form.} We favor \emph{discrete} tokens as the unified interface. Discrete tokenization integrates naturally with autoregressive modeling, enabling a single tokenized interface that unifies understanding and generation while remaining fully compatible with the LLM ecosystem~\citep{chen2025blip3o}.

\textbf{The bottleneck.} Compared to continuous features, discrete tokens typically have lower information density. For non-text modalities such as images, this creates a tight information bottleneck that can degrade multimodal understanding.

\textbf{Our approach.} We address this bottleneck with a multi-token discretization scheme that substantially increases the information capacity of discrete vision tokens, yielding understanding on par with continuous-feature VLMs while remaining fully discrete and autoregressive. For the visual modality, we use a ViT-based encoder as the Tokenizer and a diffusion-based decoder as the Detokenizer, letting each component focus on its strength: the Tokenizer aligns visual signals into the shared semantic space, the LLM handles reasoning, and the Detokenizer renders high-quality images.


\section{Architecture}
\label{sec3_arch}

\begin{figure}[t!] 
  \centering
  \includegraphics[width=0.95\textwidth]{OneBase/figures/main_archi_final.png}
    \caption{The auto-regressive inference workflow of our method, including the Kelix-Tok, and unified semantic understanding and image generation.}
    \label{mainmodel}
\end{figure}


Figure~\ref{mainmodel} shows a high-level training overview of our methods, which unify the semantic understanding\&image generation simultaneously.
%
% includes multi-token tokenizer, unified semantic understanding\&image generation, and an image diffusion de-tokenizer
In practice, our method consists by three major components:
\begin{itemize}
    \item \textbf{Kelix-Tokenizer}, which provides multiple parallel discrete tokens for each patch embedding.
    \item \textbf{Unified LLM Backbone}, to compress the high-level text and image semantic fusion.
    \item \textbf{Diffusion-based Image De-Tokenizer}, to render the high-resolution image according to last hidden states.
\end{itemize}
This section begins by introducing the discrete image tokenizer, detailing how to quantize the ViT-produced patch embeddings via a multi-token quantization mechanism.
%
We then introduce a next-block-prediction training paradigm for the unified LLM backbone, which compresses the increased vision tokens from our multi-token tokenizer into manageable blocks, thereby preserving training and inference efficiency comparable to single-token methods.
%
Finally, we describe the diffusion-based de-tokenizer that renders high-quality images conditioned on the LLM-generated image tokens.

% \begin{figure}[t!]
%   \centering
%   \includegraphics[width=0.85\textwidth]{Helix/figures/tokenizer.png}
%     \caption{The multi-token image tokenizer architecture details: for each ViT encoded continuous patch embedding, we first split it to multiple pieces, and then each piece will be quantified into a discrete index from corresponding codebook (K-means initialized), \textit{e.g.}, 8 discrete codes for a continuous patch embedding.}
%     \label{dit}
% \end{figure}





\begin{figure}[t!]
  \centering
  % \includegraphics[width=0.95\textwidth]{Helix/figures/fig3.pdf}
  \includegraphics[width=0.95\textwidth]{OneBase/figures/fig3.pdf}
    \caption{The auto-regressive training workflow of our method, including the Kelix Tokenizer, and Unified LLM model and Image DiT.}
    \label{train-pipeline}
\end{figure}





\begin{table}[t]
\footnotesize
	\centering
	\caption{A list of commonly used notation.}
	\setlength\tabcolsep{10pt}{
	\begin{tabular}{cl}
		\toprule
		\textbf{Notation} & \textbf{Description}\\
		\midrule
        \rowcolor{gray!10}
        \multicolumn{2}{l}{\textbf{Basic concept}} \\
		$I$ & An arbitrary image $I$. \\ \specialrule{0em}{2pt}{2pt}
        $P$ & The patch number of a Image $P$. \\ \specialrule{0em}{2pt}{2pt}
        $\mathbf{Z}/\mathbf{z}$ & The embedding matrix-of-all-patchs/vector-of-specific-patch. \\ \specialrule{0em}{2pt}{2pt}
        $\mathbf{W}$ & The learnable parameter matrix. \\ \specialrule{0em}{2pt}{2pt}
        \rowcolor{gray!10}
        \multicolumn{2}{l}{\textbf{VQ codebook concept}} \\
        $\mathbf{C}$ & The codebook embedding matrix. \\ \specialrule{0em}{2pt}{2pt}
        $d$ & The codebook dimension. \\ \specialrule{0em}{2pt}{2pt}
        $\mathbf{z}_q$ & The selected closest quantify embedding. \\ \specialrule{0em}{2pt}{2pt}
        \rowcolor{gray!10}
        \multicolumn{2}{l}{\textbf{Multiple VQ codebook concept}} \\
        $K$ & The codebook size, e.g., 8192. \\ \specialrule{0em}{2pt}{2pt}
        $N$ & The sub-codebook number of our Kelix, e.g., 8. \\ \specialrule{0em}{2pt}{2pt}
        $S$ & The sum of sub-codebook size of our Kelix, e.g., 8 $\times$ 8192 = 65536. \\ \specialrule{0em}{2pt}{2pt}
        \rowcolor{gray!10}
        \multicolumn{2}{l}{\textbf{Block prediction concept}} \\
        $t$ & The basic input token ID.  \\ \specialrule{0em}{2pt}{2pt}
        $B$ & The block input unit, text block is $\{t^1,\texttt{<eob>}\}$, visual block is $\{t^1,\dots,t^8,\texttt{<eob>}\}$ \\ \specialrule{0em}{2pt}{2pt}
        $\texttt{<eob>}$ & Additional special token end-of-block. \\ \specialrule{0em}{2pt}{2pt}
        $M$ & The block size, text block is 2, , visual block is 9. \\ \specialrule{0em}{2pt}{2pt}
        $\mathbf{e}$ & The input token embedding. \\ \specialrule{0em}{2pt}{2pt}
        $\mathbf{h}$ & The last hidden states. \\ \specialrule{0em}{2pt}{2pt}
		\midrule
	\bottomrule
	\end{tabular}
	\label{tab::symbol}
	}
\end{table}




% \begin{table}[t]
% 	\centering
% 	\caption{A list of commonly used notations.}
% 	\setlength\tabcolsep{10pt}{
% 	\begin{tabular}{cl}
% 		\toprule
% 		\textbf{Notation} & \textbf{Description}\\
% 		\midrule
% 		$\mathcal{U}^X$, $\mathcal{U}^Y$ & The user set. \\ \specialrule{0em}{2pt}{2pt}
% 		$\mathcal{V}^X$, $\mathcal{V}^Y$ & The item set. \\ \specialrule{0em}{2pt}{2pt}
% 		$\mathcal{U}^x$, $\mathcal{U}^y$ & The non-overlapping (cold-start) user set. \\ \specialrule{0em}{2pt}{2pt}
% 		$\mathcal{U}^o$ & The overlapping user set across domains. \\ \specialrule{0em}{2pt}{2pt}
% 		$\bm{A}^X$, $\bm{A}^Y$ & The binary adjacent matrix for user-item interaction. \\ \specialrule{0em}{2pt}{2pt}
% 		\midrule
% 		$\bm{Z}_u^X$, $\bm{Z}_u^Y$ & User representations. \\ \specialrule{0em}{2pt}{2pt}
% % 		$\bm{Z}_v^X$, $\bm{Z}_v^Y$ & Items' representations. \\ \specialrule{0em}{2pt}{2pt}	
%         $\bm{Z}_v^X$, $\bm{Z}_v^Y$ & Item representations. \\ \specialrule{0em}{2pt}{2pt}
% % 		$\bm{Z}_u^{xo}$, $\bm{Z}_u^{yo}$ & Overlapping users' representations.\\ \specialrule{0em}{2pt}{2pt}
% 		$\bm{Z}_u^{xo}$, $\bm{Z}_u^{yo}$ & Overlapping user representations.\\ \specialrule{0em}{2pt}{2pt}
% % 		$\bm{Z}_u^{x}$, $\bm{Z}_u^{y}$ & Non-overlapping (cold-start) users' representations.\\ \specialrule{0em}{2pt}{2pt}
%         $\bm{Z}_u^{x}$, $\bm{Z}_u^{y}$ & Non-overlapping (cold-start) user representations.\\ \specialrule{0em}{2pt}{2pt}
% 		$\mathbf{X}$, $\mathbf{Y}$ & The user-item interaction information. \\ \specialrule{0em}{2pt}{2pt}
% % 		& i.e., the observed user-item behaviors. \\ \specialrule{0em}{2pt}{2pt}
% % 		$\mathbf{X}^u$, $\mathbf{Y}^u$ & The observed user information. \\ \specialrule{0em}{2pt}{2pt}
%         $\mathbf{X}^u$, $\mathbf{Y}^u$ & The user information (users' features and structure). \\ \specialrule{0em}{2pt}{2pt}
% % 		& i.e., the user feature and the related user information. \\ \specialrule{0em}{2pt}{2pt}
% 		$\mathbf{X}^v$, $\mathbf{Y}^v$ & The item information (items' features and structure). \\ \specialrule{0em}{2pt}{2pt}
% % 		& i.e., the item feature and the related item information. \\ \specialrule{0em}{2pt}{2pt}
% 		\midrule
% 		$\mathcal{L}_{o2X}$, $\mathcal{L}_{o2Y}$ & Cross-domain information bottleneck regularizer. \\ \specialrule{0em}{2pt}{2pt}
% 		$\mathcal{L}_{x2X}$, $\mathcal{L}_{y2Y}$ & In-domain information bottleneck regularizer. \\ \specialrule{0em}{2pt}{2pt}
% 		$\mathcal{L}_{con}$ & Contrastive information regularizer. \\ \specialrule{0em}{2pt}{2pt}
% 	\bottomrule
% 	\end{tabular}
% 	\label{tab::symbol}
% 	}
% \end{table}


\subsection{Discrete Image Tokenizer}
\label{sec:discrete_image_tokenizer}
Recent works on discrete image tokenizers, such as TAToK~\citep{kim2025democratizing}, BLIP3o-Next~\citep{chen2025blip3o}, Emu3~\citep{wang2024emu3}, and DualToken~\citep{song2025dualtoken}, X-omni~\cite{geng2025xomnireinforcementlearningmakes}, have demonstrated strong image generation capabilities.
%
However, these methods generally underperform VLMs that operate on continuous vision features in multi-modal understanding benchmarks (i.e., OCRbench, most recent VLMs achieve 85+, while the stongest unified model with discrete tokens only reach 70.). We attribute this gap to the inherent information bottleneck of single-token quantization. Consider, for example, a 1024-dimensional continuous embedding quantized into a single discrete token with a vocabulary size of 65{,}536: the resulting coding capacity is merely $\log_2 65536 = 16$ bits, whereas the original floating-point embedding occupies approximately 2\,KB. Although continuous embeddings inevitably contain some redundancy, their effective information capacity still far exceeds what 16 bits can represent.
%
To bridge this gap, we propose quantizing each continuous patch embedding into \emph{multiple} parallel discrete tokens, thereby significantly expanding the representational capacity. After quantization, we fuse these tokens back into a single composite representation before feeding it to the LLM decoder, which avoids the context length explosion that would otherwise result from the increased token count.
%
In the following, we describe the core design of Kelix-Tok in detail. An overview of the framework is presented on the left of Figure~\ref{train-pipeline}.


\textbf{Vector Quantization.}
Vector Quantization (VQ)~\citep{van2017neural} discretizes continuous representations into a finite set of learnable codes. Given an image $I$ with $P$ patches, a vision encoder produces a continuous feature matrix $\mathbf{Z}_{I} = [\mathbf{z}_{1}, \mathbf{z}_{2}, \cdots, \mathbf{z}_{P}] = \mathrm{ViT}(I) \in \mathbb{R}^{P \times d}$, where $d$ is the embedding dimension. Standard VQ maps each patch feature $\mathbf{z}$ to its nearest entry in a codebook $\mathbf{C} = \{\mathbf{c}_{1}, \mathbf{c}_{2}, \ldots, \mathbf{c}_{K}\}$:
%
\begin{equation}
    \mathbf{z}_{q} = \mathbf{c}_{\texttt{Index}}, \quad \texttt{Index} = \underset{\mathbf{c} \in \mathbf{C}}{\arg\min} \big\| \mathbf{z} - \mathbf{c} \big\|^{2}
\end{equation}
%
where $K$ is the codebook size. As discussed above, this single-token quantization creates a severe information bottleneck that the coding capacity of a single discrete index is far below the effective information carried by the original continuous embedding.

\textbf{Multi-Token Quantization.}
To address this bottleneck, we decompose each continuous patch feature into $N$ parallel subspaces via dedicated projection matrices, and quantize each subspace independently with its own sub-codebook. The combinatorial coding capacity thus grows exponentially with $N$, substantially narrowing the gap with continuous representations.
%
To keep the context length unchanged for the LLM encoder, we apply \textbf{sum pooling
}to aggregate the $N$ quantized sub-vectors back into a single composite token on the encoder side, while the decoder retains the full $N$-way discrete codes for autoregressive prediction. This asymmetric design balances representational expressivity with computational efficiency.

\textbf{Codebook Initialization.}
Rather than random initialization, we leverage a pre-trained Keye-VL (NaViT) vision encoder to extract visual patch embedding from a large image corpus (more than 10,000,000 images), and apply K-means clustering to obtain $S = 65{,}536$ high-quality cluster centers. These centers are then uniformly partitioned into $N$ non-overlapping sub-codebooks, each containing $S / N$ entries, ensuring complementary semantic coverage across subspaces. During training, we adopt the SimVQ~\citep{zhu2025addressing} strategy: the cluster centers are frozen and only the per-subspace SimVQ projection matrices $\mathbf{W}_i$ are optimized, which maximizes codebook utilization and effectively mitigates codebook collapse.

\textbf{Architecture and Training Objective.}
Kelix-Tok comprises two core components: (1)~a NaViT-style vision encoder, initialized from Keye-VL and supporting native-resolution input, and (2)~the multi-token quantization layer described above, which decomposes each continuous patch feature into $N$ parallel discrete codes.
%
To align the resulting discrete visual representations with the LLM text space, we attach a pre-trained lightweight language model (600M parameters) as an auxiliary decoder during training. The overall objective jointly optimizes a next-token prediction loss and a quantization loss:
%
\begin{equation}
\mathcal{L} = \mathcal{L}_{\text{NTP}} + \mathcal{L}_{\text{VQ}}
\end{equation}
%
where $\mathcal{L}_{\text{NTP}}$ is the standard cross-entropy loss that supervises the auxiliary decoder, enforcing semantic alignment between discrete visual tokens and the language space. The quantization loss $\mathcal{L}_{\text{VQ}}$ is defined as:
%
\begin{equation}
\mathcal{L}_{\text{VQ}} = \dfrac{1}{N}\sum_{i=1}^{N} \left( \underbrace{\big\| \text{sg}[\mathbf{z}^{\left(i\right)}] - \mathbf{z}_q^{\left(i\right)} \big\|_2^2}_{\text{Codebook Loss}} + \underbrace{\beta \big\| \mathbf{z}^{\left(i\right)} - \text{sg}[\mathbf{z}_q^{\left(i\right)}] \big\|_2^2 }_{\text{Commitment Loss}}\right)
\end{equation}
%
where $\mathbf{z}^{(i)}$ and $\mathbf{z}_q^{(i)}$ denote the continuous and quantized features from the $i$-th subspace, respectively, $\text{sg}[\cdot]$ is the stop-gradient operator, and $\beta$ controls the commitment weight. Intuitively, the codebook loss pulls each codebook entry toward the corresponding encoder output, while the commitment loss encourages the encoder to produce features that stay close to their quantized counterparts.

\textbf{Training Strategy.}
We adopt a staged unfreezing schedule: model parameters are gradually unlocked and trained on a diverse collection of understanding-oriented datasets spanning commonsense reasoning, inference, and text-intensive tasks. This curriculum strengthens the tokenizer's semantic comprehension while preserving the perceptual capabilities inherited from NaViT. Once training is complete, the auxiliary language model decoder is discarded; only the vision tokenizer is retained and integrated into the downstream larger unified LLM and image de-tokenizer pipeline.



\subsection{Unified Model Architecture}

As discussed in Section~\ref{sec:discrete_image_tokenizer}, our multi-token quantization produces $N$ discrete codes per patch, substantially increasing the number of vision tokens compared to the original ViT output. To handle this without inflating the context length, we adopt a unified autoregressive Transformer architecture built around a \textbf{Next-Block Prediction (NBP)} paradigm.

\textbf{Next-Block Prediction (NBP).}
We organize multimodal sequences into a series of blocks $\{B_1, B_2, \dots\}$, where each block $B_i$ represents a coherent semantic unit terminated by an end-of-block token (\texttt{<eob>}). The block size is modality-dependent: 
%
\begin{itemize}[noitemsep, topsep=2pt]
    \item \myorange{\textbf{Text block ($\{t^1, \texttt{<eob>}\}$)}}: contains 1 text token plus \texttt{<eob>} ($|B|=2$), preserving token-level granularity consistent with standard text LLMs.
    \item \myorange{\textbf{Visual block ($\{t^1, \dots, t^N, \texttt{<eob>}\}$)}}: contains $N$ visual tokens plus \texttt{<eob>} ($|B|=N{+}1$), compressing the expanded discrete codes of each patch into a single prediction step while retaining local spatial structure.
\end{itemize}
At each autoregressive step $i$, the LLM processes block-level representations $\{B_1, \dots, B_i\}$ and predicts the entire next block $B_{i+1}$, effectively reducing the sequence length seen by the backbone by a factor of $N$ for visual data.

% This paradigm elevates the LLM’s processing to a higher semantic granularity: it reduces redundant token-level computation for visual data (a key limitation of NTP in multimodal scenarios), preserves fine-grained details via intra-block modeling, and maintains full compatibility with the autoregressive Transformer backbone. Critically, NBP ensures seamless unification of text and vision processing without modifying the core architecture, laying the foundation for efficient unified understanding and generation while inheriting NTP’s scaling advantages.

\textbf{Block-wise Input and Output Modules.}
To interface the LLM backbone with block-based sequences, we introduce two lightweight components: the \textit{Block Encoder} and the \textit{Block Decoder}. \myorange{For notational brevity, we denote each block as $B_i = \{t_{i}^1, t_{i}^2, \dots, t_{i}^M\}$, where $M = 2$ for text blocks and $M = N{+}1$ for visual blocks.}

\textit{Block Encoder (Embedding Aggregation):}
The \textit{Block Encoder} maps each block into a single fixed-dimensional representation for the LLM. Given the token embeddings $\{\mathbf{e}_{i}^1, \mathbf{e}_{i}^2, \dots, \mathbf{e}_{i}^M\}$ of block $B_i$, we compute the element-wise sum \myred{excluding the end-of-block embedding $\mathbf{e}_i^M$}:
\begin{equation}
\label{eq:ea_1}
    \mathbf{E}_i = \sum_{j=1}^{M-1} \mathbf{e}_{i}^j.
\end{equation}
This yields a modality-agnostic input to the LLM. The forward pass is then:
\begin{equation}
\label{eq:ea}
    \{\mathbf{h}_1,\dots,\mathbf{h}_{i}\} = \texttt{LLM}\left(\{\mathbf{E}_1,\dots,\mathbf{E}_{i}\}\right)
\end{equation}
where $\mathbf{h}_i$ denotes the last-layer hidden state at position $i$.

\textit{Block Decoder (Autoregressive Reconstruction):}
Given the hidden state $\mathbf{h}_i$, which encodes all preceding blocks $B_{\leq i}$, a lightweight Transformer decoder autoregressively reconstructs the next block:
\begin{equation}
    B_{i+1} = \texttt{BlockDecoder}\left(\mathbf{h}_i\right).
    \label{blockdecoder}
\end{equation}
The corresponding training loss factorizes as:
%
\begin{align*}
\mathcal{L}_{\text{nbp}}(B_{i+1}) &= -\log P\left(B_{i+1}\big|B_{\leq i}\right) \\
                  &= - \sum_{j=1}^{M} \log P\left(t^{j}_{i+1}\big|\mathbf{h}_{i},\, t^{<j}_{i+1}\right)
\end{align*}
%
where $\mathcal{L}_{\text{nbp}}$ serves as the training objective for the unified model.

\textbf{Implementation Details.}
We use Qwen3-8B~\citep{yang2025qwen3} as the LLM backbone. To accommodate the discrete visual codes from Kelix-Tok, we expand the vocabulary by $S = 65{,}536$ additional entries and augment the LM head accordingly. With the NBP paradigm, the effective context length for visual data is reduced by a factor of $N$, yielding training and inference efficiency comparable to single-token approaches. 

% \begin{figure}[t!]
%   \centering
%   \includegraphics[width=16cm,height=5.3cm]{Helix/figures/DiT.png}
%     \caption{The diffusion based image de-tokenizer training.}
%     \label{dit}
% \end{figure}





\subsection{Diffusion-based Image De-Tokenizer}
The unified model described above produces two types of output: text tokens and image tokens. Text tokens can be directly mapped back to words via the standard text tokenizer. Image tokens, however, encode high-level semantic information that cannot be converted into pixels through a simple deterministic decoder (e.g., a VAE decoder~\citep{yang2024cogvideox}).
%
To render high-quality images from these semantic representations, we employ a latent diffusion model as the image de-tokenizer. Specifically, given a generation prompt, the Kelix LLM autoregressively produces a sequence of last hidden states $\{\mathbf{h}_*\}$ corresponding to the image blocks. We use these continuous hidden states as the semantic condition to guide a Diffusion Transformer (DiT). Following the flow-matching paradigm, we fine-tune a customized DiT (based on SANA-DiT~\citep{xie2024sana}) to denoise in the latent space of a pre-trained VAE, conditioned on $\{\mathbf{h}_*\}$. An overview is shown on the right of Figure~\ref{train-pipeline}.

\begin{table}[t]
    \caption{
        Comprehensive comparison on a curated set of multimodal benchmarks. The release dates denote the time of arXiv submission or official blog release.
    }
    \centering
    \small
    \renewcommand{\arraystretch}{1.2}
    \resizebox{\textwidth}{!}{
    % 去掉了所有的 |
    \begin{tabular}{l c  ccc  ccc  ccc}
    \toprule
    \multirow{2}{*}{\textbf{Model}} & \multirow{2}{*}{\textbf{Date}} & \multicolumn{3}{c}{\textbf{General Benchmarks}} & \multicolumn{3}{c}{\textbf{Knowledge Benchmarks}} & \multicolumn{3}{c}{\textbf{Text-rich Benchmarks}} \\
    
    % 使用 (lr) 来让横线产生间隙，从而在视觉上起到分隔列的作用
    \cmidrule(lr){3-5} \cmidrule(lr){6-8} \cmidrule(lr){9-11}
    
    & & {SEED-Bench} & {RealWorldQA} & {MMBench$_\text{EN}$} & {AI2D$_\text{test}$} & {MMMU} & {MathVista} & {ChartQA} & {TextVQA} & {OCRBench} \\
    \midrule
    
    \rowcolor{gray!10}
    \multicolumn{11}{l}{\textbf{Understanding-only Models}} \\
    LLaVA-OV (7B)         & 2024.08   & 76.7 & 69.9 & -- & 81.4 & 48.8 & -- & 80.9 & -- & 62.2 \\
    BLIP-3 (4B)           & 2024.08   & 72.2 & 60.5 & -- & --   & 41.1 & 39.6 & -- & 71.0 & -- \\
    InternVL3.5 (4B)      & 2025.08   & --   & 66.3 & -- & 82.6 & 66.6 & 77.1 & 86.0 & 77.9 & 82.2 \\
    Qwen2.5VL (3B)        & 2025.02   & --   & 65.4 & 76.4 & 81.6 & 53.1 & 62.3 & 84.0 & 79.7 & 79.7 \\
    Qwen2.5VL (7B)        & 2025.02   & --   & 68.5 & 82.6 & 83.9 & 58.6 & 68.2 & 87.3 & 84.9 & 86.4 \\
    \midrule

    \rowcolor{gray!10}
    \multicolumn{11}{l}{\textbf{Unified Multimodal LLMs}} \\
    
    \multicolumn{11}{c}{\textbf{\textit{Continuous Tokenization}}} \\
    Bagel (14B)           & 2025.05   & 78.5 & 72.8 & 85.0 & 89.2 & 55.3 & 73.1 & 78.5 & 80.0 & 78.3 \\
    Ming-Univision (16B)  & 2025.10   & --   & --   & --   & 82.8   & 40.3   & 66.6   & --   & --   & 72.4   \\
    Blip-3o (4B)          & 2025.05   & 73.8 & 60.4 & 78.6 & --   & 46.6 & --   & --   & 78.0 & --   \\
    TUNA (7B)             & 2025.12   & 74.7 & 66.1 & --   & 79.3 & 49.8 & --   & 85.8 & --   & 74.3 \\

    \multicolumn{11}{c}{\textbf{\textit{Hybrid Tokenization}}} \\
    Manzano (30B)         & 2025.09   & 76.0 & 70.1 & 83.4 & 86.0 & 57.8 & 73.3 & 89.0 & 84.3 & 86.3 \\
    VQ-RAE (7B)  & 2025.11   & 77.0 & -- & 85.1 & 84.8 & 61.6 & -- & -- & 80.6 & -- \\
    \rowcolor{blue!10}
    \multicolumn{11}{c}{\textbf{\textit{Discrete Tokenization}}} \\
    \rowcolor{blue!10}
    Janus-Pro (1.5B)      & 2025.01   & 68.3 & 52.6 & 75.5 & 64.5 & 36.3 & 36.5 & 23.4 & 41.9 & 48.7 \\
    \rowcolor{blue!10}
    Janus-Pro (7B)        & 2025.01   & 72.1 & 58.0 & \underline{79.2} & 71.3 & 41.0 & 42.5 & 25.8 & 45.6 & 59.0 \\
    \rowcolor{blue!10}
    Emu3 (8B)             & 2024.09   & 68.2 & 57.4 & 58.5 & 70.0 & 31.6 & 47.6 & --   & 64.7 & 68.7 \\
    \rowcolor{blue!10}
    X-Omni (7B)           & 2025.07   & \underline{74.3} & \underline{62.6} & 74.8 & 76.8 & 47.2 & \underline{54.1} & \underline{81.5} & \underline{77.4} & \underline{70.4} \\
    \rowcolor{blue!10}
    Show-O2 (7B)          & 2025.06   & 69.8   & --   & --   & \underline{78.6}   & \underline{48.9}   & --   & 66.9   & 71.5   & --   \\
    \rowcolor{blue!10}
    Unitok (7B)           & 2025.02   & --   & --   & --   & --   & --   & --   & --   & 51.6   & --   \\
    \rowcolor{blue!10}
    TokenFlow-XL (14B)       & 2024.12   & 72.6   & 56.6   & --   & 75.8   & 43.2   & --   & --   & 62.3   & --   \\
    \rowcolor{blue!10}
    Nextflow (7B)         & 2026.01   & --   & --   & --   & --   & 37.1   & --   & 57.7   & 58.9   & 55.1   \\
    \rowcolor{blue!10}
    Kelix (8B)          &           & \textbf{76.0} & \textbf{72.1} & \textbf{80.2} & \textbf{82.4} & \textbf{54.1} & \textbf{76.5} & \textbf{83.0} & \textbf{81.4} & \textbf{86.7} \\
    \midrule
    
    \rowcolor{gray!10}
    \multicolumn{11}{l}{\textbf{Proprietary Models}} \\
    GPT-4o                & 2024.05   & 77.1 & 75.4 & --   & 84.6 & 69.2 & 61.3 & 85.7 & --   & 73.6 \\
    Gemini-2.5-Pro        & 2024.12   & --   & 78.0 & 86.3 & 89.5 & 81.7 & 82.7 & 83.3 & 76.8 & 86.2 \\
    
    \bottomrule
    \end{tabular}
    }
\label{tab:full_comprehensive_evaluation_clean}
For a fair comparison, baselines highlighted with \colorbox{blue!10}{blue shading} adopt the same discrete tokenization paradigm and comparable model scale as Kelix. \textbf{Best} and \underline{runner-up} results among these comparable models are marked in \textbf{bold} and \underline{underlined}, respectively.
\end{table}

% \begin{table}[t]
%     \caption{
%         Comprehensive comparison on a curated set of multimodal benchmarks. The release dates denote the time of arXiv submission or official blog release.
%     }
%     \centering
%     \small
%     \renewcommand{\arraystretch}{1.2}
%     \resizebox{\textwidth}{!}{
%     \begin{tabular}{l c | ccc | ccc | ccc}
%     \toprule
%     % --- 第一层表头 ---
%     \multirow{2}{*}{\textbf{Model}} & \multirow{2}{*}{\textbf{Date}} & \multicolumn{3}{c|}{\textbf{General Benchmarks}} & \multicolumn{3}{c|}{\textbf{Knowledge Benchmarks}} & \multicolumn{3}{c}{\textbf{Text-rich Benchmarks}} \\
    
%     \cmidrule(lr){3-5} \cmidrule(lr){6-8} \cmidrule(lr){9-11}
    
%     % --- 第二层表头 ---
%     & & {SEED-Bench} & {RealWorldQA} & {MMBench$_\text{EN}$} & {AI2D$_\text{test}$} & {MMMU} & {MathVista} & {ChartQA} & {TextVQA} & {OCRBench} \\
%     \midrule
    
%     \rowcolor{gray!10}
%     \multicolumn{11}{l}{\textbf{Understanding-only Models}} \\
%     LLaVA-OV (7B)         & 2024.08   & 76.7 & 69.9 & -- & 81.4 & 48.8 & -- & 80.9 & -- & 62.2 \\
%     BLIP-3 (4B)           & 2024.08   & 72.2 & 60.5 & -- & --   & 41.1 & 39.6 & -- & 71.0 & -- \\
%     InternVL3.5 (4B)      & 2025.08   & --   & 66.3 & -- & 82.6 & 66.6 & 77.1 & 86.0 & 77.9 & 82.2 \\
%     Qwen2.5VL (3B)        & 2025.02   & --   & 65.4 & 76.4 & 81.6 & 53.1 & 62.3 & 84.0 & 79.7 & 79.7 \\
%     Qwen2.5VL (7B)        & 2025.02   & --   & 68.5 & 82.6 & 83.9 & 58.6 & 68.2 & 87.3 & 84.9 & 86.4 \\
%     \midrule

%     \rowcolor{gray!10}
%     \multicolumn{11}{l}{\textbf{Unified Multimodal LLMs}} \\
    
%     % --- Continuous Sub-section ---
%     \multicolumn{11}{c}{\textbf{\textit{Continuous Tokenization}}} \\
%     Bagel (14B)           & 2025.05   & 78.5 & 72.8 & 85.0 & 89.2 & 55.3 & 73.1 & 78.5 & 80.0 & 78.3 \\
%     Ming-Univision (16B)  & 2025.10   & --   & --   & --   & 82.8   & 40.3   & 66.6   & --   & --   & 72.4   \\
%     Blip-3o (4B)          & 2025.05   & 73.8 & 60.4 & 78.6 & --   & 46.6 & --   & --   & 78.0 & --   \\
%     TUNA (7B)             & 2025.12   & 74.7 & 66.1 & --   & 79.3 & 49.8 & --   & 85.8 & --   & 74.3 \\

%     % --- Hybrid Sub-section ---
%     \multicolumn{11}{c}{\textbf{\textit{Hybrid Tokenization}}} \\
%     % Manzano (3B)          & 2025.09   & 74.3 & 65.1 & 78.1 & 82.2 & 51.4 & 69.8 & 88.2 & 80.1 & 85.7 \\
%     Manzano (30B)         & 2025.09   & 76.0 & 70.1 & 83.4 & 86.0 & 57.8 & 73.3 & 89.0 & 84.3 & 86.3 \\

%     % --- Discrete Sub-section ---
%     \multicolumn{11}{c}{\textbf{\textit{Discrete Tokenization}}} \\
%     Janus-Pro (1.5B)      & 2025.01   & 68.3 & 52.6 & 75.5 & 64.5 & 36.3 & 36.5 & 23.4 & 41.9 & 48.7 \\
%     Janus-Pro (7B)        & 2025.01   & 72.1 & 58.0 & 79.2 & 71.3 & 41.0 & 42.5 & 25.8 & 45.6 & 59.0 \\
%     Emu3 (8B)             & 2024.09   & 68.2 & 57.4 & 58.5 & 70.0 & 31.6 & 47.6 & --   & 64.7 & 68.7 \\
%     X-Omni (7B)           & 2025.07   & 74.3 & 62.6 & 74.8 & 76.8 & 47.2 & 54.1 & 81.5 & 77.4 & 70.4 \\
%     Show-O2 (7B)          & 2025.06   & 69.8   & --   & --   & 78.6   & 48.9   & --   & 66.9   & 71.5   & --   \\
%     Unitok (7B)           & 2025.02   & --   & --   & --   & --   & --   & --   & --   & 51.6   & --   \\
%     TokenFlow-XL (14B)       & 2024.12   & 72.6   & 56.6   & --   & 75.8   & 43.2   & --   & --   & 62.3   & --   \\
%     Nextflow (7B)         & 2026.01   & --   & --   & --   & --   & 37.1   & --   & 57.7   & 58.9   & 55.1   \\
%     Nextflow-RL (7B)      & 2026.01   & --   & --   & --   & --   & --   & --   & --   & --   & --   \\
%     \rowcolor{blue!10}
%     Kelix (8B)          &           & 76.0 & 72.1 & 80.2 & 82.4 & 54.1 & 76.5 & 83.0 & 81.4 & 86.7 \\
%     \midrule
    
%     \rowcolor{gray!10}
%     \multicolumn{11}{l}{\textbf{Proprietary Models}} \\
%     GPT-4o                & 2024.05   & 77.1 & 75.4 & --   & 84.6 & 69.2 & 61.3 & 85.7 & --   & 73.6 \\
%     Gemini-2.5-Pro        & 2024.12   & --   & 78.0 & 86.3 & 89.5 & 81.7 & 82.7 & 83.3 & 76.8 & 86.2 \\
    
%     \bottomrule
%     \end{tabular}
%     }
%     \label{tab:full_comprehensive_evaluation_clean}
% \end{table}



% \begin{table}[t]
%     \caption{
%         Comprehensive comparison of image generation performance on GenEval and WISE benchmarks. The release dates denote the time of arXiv submission or official release. \textbf{For Kelix, we report performance using the default resolution of $1024\times1024$.}
%     }
%     \centering
%     \small
%     \renewcommand{\arraystretch}{1.2}
%     \resizebox{\textwidth}{!}{
%     \begin{tabular}{l c | ccccccc | ccccccc}
%     \toprule
%     \multirow{2}{*}{\textbf{Model}} & \multirow{2}{*}{\textbf{Date}} & \multicolumn{7}{c|}{\textbf{GenEval Benchmark}} & \multicolumn{7}{c}{\textbf{WISE Benchmark}} \\
%     \cmidrule(lr){3-9} \cmidrule(lr){10-16}
%     & & Single & Two & Count & Color & Pos. & C-Attr & \textbf{Overall} & Cult. & Time & Space & Bio. & Phys. & Chem. & \textbf{Overall} \\
%     \midrule
    
%     \rowcolor{gray!10}
%     \multicolumn{16}{l}{\textbf{Dedicated T2I Models}} \\
%     SDXL (3.5B)           & 2023.07   & 98.0 & 74.0 & 39.0 & 85.0 & 15.0 & 23.0 & 55.0 & 43.0 & 48.0 & 47.0 & 44.0 & 45.0 & 27.0 & 43.0 \\
%     DALL-E 3              & 2023.10   & 96.0 & 87.0 & 47.0 & 83.0 & 43.0 & 45.0 & 67.0 & --   & --   & --   & --   & --   & --   & --   \\
%     SD3 Medium (2B)       & 2024.06   & 99.0 & 94.0 & 72.0 & 89.0 & 33.0 & 60.0 & 74.0 & 43.0 & 50.0 & 52.0 & 41.0 & 53.0 & 33.0 & 45.0 \\
%     FLUX.1-dev (12B)      & 2024.08   & 98.0 & 93.0 & 75.0 & 93.0 & 68.0 & 65.0 & 82.0 & 48.0 & 58.0 & 62.0 & 42.0 & 51.0 & 35.0 & 50.0 \\
%     \midrule

%     \rowcolor{gray!10}
%     \multicolumn{16}{l}{\textbf{LLM \& Diffusion Conjunction}} \\
%     MetaQuery-XL (7B)     & 2025.05   & --   & --   & --   & --   & --   & --   & 80.0 & 56.0 & 55.0 & 62.0 & 49.0 & 63.0 & 41.0 & 55.0 \\
%     OmniGen-2 (7B)        & 2025.06   & 100.0 & 95.0 & 64.0 & 88.0 & 55.0 & 76.0 & 80.0 & --   & --   & --   & --   & --   & --   & --   \\
%     % Qwen-Image (27B)      & 2025.08   & 99.0 & 92.0 & 89.0 & 88.0 & 76.0 & 77.0 & 87.0 & 67.0 & 67.0 & 80.0 & 62.0 & 79.0 & 41.0 & 67.0 \\
%     Qwen-Image (27B)      & 2025.08   & 99.0 & 92.0 & 89.0 & 88.0 & 76.0 & 77.0 & 87.0 & -- & -- & -- & -- & -- & -- & -- \\
%     \midrule

%     \rowcolor{gray!10}
%     \multicolumn{16}{l}{\textbf{Unified Multimodal LLMs}} \\
    
%     % --- Continuous Sub-section (Centered) ---
%     \multicolumn{16}{c}{\textbf{\textit{Continuous Tokenization}}} \\
%     Bagel (14B)           & 2025.05   & 99.0 & 94.0 & 81.0 & 88.0 & 64.0 & 63.0 & 82.0 & 44.0 & 55.0 & 68.0 & 44.0 & 60.0 & 39.0 & 52.0 \\
%     Ming-Univision (16B)  & 2025.10   & 100.0   & 93.0   & 59.0   & 93.0   & 92.0   & 70.0   & 85.0 & --   & --   & --   & --   & --   & --   & --   \\
%     Blip-3o (4B)          & 2025.05   & --   & --   & --   & --   & --   & --   & 81.0 & --   & --   & --   & --   & --   & --   & 50.0 \\
%     TUNA (7B)             & 2025.12   & 100   & 97   & 81   & 91   & 88   & 83   & 90.0 & --   & --   & --   & --   & --   & --   & --   \\

%     % --- Hybrid Sub-section (Centered) ---
%     \multicolumn{16}{c}{\textbf{\textit{Hybrid Tokenization}}} \\
%     % Manzano (3B)          & 2025.09   & 98.0 & 91.0 & 82.0 & 71.0 & 78.0 & 71.0 & 85.0 & 42.0 & 51.0 & 59.0 & 45.0 & 51.0 & 32.0 & 46.0 \\
%     Manzano (30B)         & 2025.09   & 100.0 & 91.0 & 83.0 & 87.0 & 84.0 & 65.0 & 85.0 & 58.0 & 50.0 & 65.0 & 50.0 & 55.0 & 32.0 & 54.0 \\

%     % --- Discrete Sub-section (Centered) ---
%     \multicolumn{16}{c}{\textbf{\textit{Discrete Tokenization}}} \\
%     Janus-Pro (7B)        & 2025.01   & 99.0 & 89.0 & 59.0 & 90.0 & 79.0 & 66.0 & 80.0 & 30.0 & 37.0 & 49.0 & 36.0 & 42.0 & 26.0 & 35.0 \\
%     Emu3 (8B)             & 2024.09   & 99   & 81   & 42   & 80   & 49   & 45   & 66.0 & --   & --   & --   & --   & --   & --   & 39.0 \\
%     X-Omni (7B)           & 2025.07   & 98.0 & 95.0 & 75.0 & 91.0 & 71.0 & 68.0 & 83.0 & --   & --   & --   & --   & --   & --   & --   \\
%     Show-O2 (7B)          & 2025.06   & 100.0   & 87.0   & 58.0   & 92.0   & 52.0   & 62.0 & 76.0 & --   & --   & --   & --   & --   & --   & --   \\
%     Unitok (7B)           & 2025.02   & 99.0   & 71.0   & 36.0   & 79.0   & 26.0   & 45.0   & 59.0 & --   & --   & --   & --   & --   & --   & --   \\
%     TokenFlow (13B)       & 2024.12   & 97.0   & 66.0   & 40.0   & 84.0   & 17.0   & 26.0   & 55.0 & --   & --   & --   & --   & --   & --   & --   \\
%     Nextflow (7B)         & 2026.01   & 98.0   & 92.0   & 73.0   & 90.0   & 77.0   & 69.0   & 83.0 & 62.0   & 60.0   & 70.0   & 54.0   & 58.0   & 38.0   & 59.0 \\
%     Nextflow-RL (7B)      & 2026.01   &  100.0  & 92.0   & 75.0   & 90.0   & 76.0   & 70.0   & 84.0 & 63.0   & 63.0   & 77.0   & 58.0   & 67.0   & 39.0   & 62.0 \\
%     %\rowcolor{blue!10}
%     %Kelix-TC-DiT (8B)          &           & 99.0 & 93.0 & 84.0 & 85.0 & 84.0 & 70.0 & 85.8 & 53.0 & 58.0 & 66.0 & 43.0 & 54.0 & 32.0 & 53.0 \\
%     \rowcolor{blue!10}
%     % Kelix (8B)          &           & 99.0 & 94.0 & 76.0 & 91.0 & 78.0 & 81.0 & 86.6 & 58.0 & 62.0 & 68.0 & 46.0 & 61.0 & 35.0 & 57.0 \\
%     Kelix (8B) & & 100.0 &	93.0 &		81.0 &		87.0 &		85.0 &		79.0 &		87.6 &		58.0	 &	62.0	 &	68.0	 &	47.0	 &	61.0	 &	35.0	 &	57.0 \\
%     \midrule
    
%     \rowcolor{gray!10}
%     \multicolumn{16}{l}{\textbf{Proprietary Models}} \\
%     GPT-4o                & 2024.05   & 99.0 & 92.0 & 85.0 & 92.0 & 75.0 & 61.0 & 84.0 & 81.0 & 71.0 & 89.0 & 83.0 & 79.0 & 74.0 & 80.0 \\
    
%     \bottomrule
%     \end{tabular}
%     }
%     \label{tab:generation_results}
% \end{table}

% \begin{table}[t!]
%     \caption{
%         Comparison of image generation performance on the DPG-Bench benchmark. All scores are reported on a 100-point scale. The release dates denote the time of arXiv submission or official release.
%     }
%     \centering
%     \small
%     \renewcommand{\arraystretch}{1.2}
%     \resizebox{0.9\textwidth}{!}{
%     % 8 columns: Model(1) + Date(1) + Scores(5) + Overall(1)
%     \begin{tabular}{l c | ccccc | c}
%     \toprule
%     \textbf{Method} & \textbf{Date} & \textbf{Global} & \textbf{Entity} & \textbf{Attribute} & \textbf{Relation} & \textbf{Other} & \textbf{Overall$\uparrow$} \\
%     \midrule
    
%     \rowcolor{gray!10}
%     \multicolumn{8}{l}{\textbf{Dedicated T2I Models}} \\
%     SDXL (2.6B)           & 2023.07 & 83.3 & 82.4 & 80.9 & 86.8 & 80.4 & 74.7 \\
%     DALL-E 3              & 2023.10 & 91.0 & 89.6 & 88.4 & 90.6 & 89.8 & 83.5 \\
%     SD3 Medium (2B)       & 2024.06 & 87.9 & 91.0 & 88.8 & 80.7 & 88.7 & 84.1 \\
%     FLUX.1-dev (12B)      & 2024.08 & --   & --   & --   & --   & --   & --   \\
%     \midrule

%     \rowcolor{gray!10}
%     \multicolumn{8}{l}{\textbf{LLM \& Diffusion Conjunction}} \\
%     MetaQuery-XL (7B)     & 2025.04 & --   & --   & --   & --   & --   & 82.05  \\
%     OmniGen-2 (7B)        & 2025.06 & 88.81   & 88.83   & 90.18   & 89.37   & 90.27   & 83.57   \\
%     Qwen-Image (27B)      & 2025.08 & 91.32   & 91.56   & 92.02   & 94.31   & 92.73   & 88.32   \\
%     \midrule

%     \rowcolor{gray!10}
%     \multicolumn{8}{l}{\textbf{Unified Multimodal LLMs}} \\
    
%     % --- Continuous Sub-section (Span 8) ---
%     \multicolumn{8}{c}{\textbf{\textit{Continuous Tokenization}}} \\
%     Bagel (14B)           & 2025.05 & --   & --   & --   & --   & --   & --   \\
%     Ming-Univision (16B)  & 2025.10 & --   & --   & --   & --   & --   & 82.1 \\
%     Blip-3o (4B)          & 2025.05 & --   & --   & --   & --   & --   & 79.4 \\
%     TUNA (7B)             & 2025.12 & 90.42   & 91.58   & 90.94   & 91.87   & 90.73   & 86.8 \\

%     % --- Hybrid Sub-section (Span 8) ---
%     \multicolumn{8}{c}{\textbf{\textit{Hybrid Tokenization}}} \\
%     % Manzano (3B)          & 2025.09 & --   & --   & --   & --   & --   & --   \\
%     Manzano (30B)         & 2025.09 & --   & --   & --   & --   & --   & --   \\

%     % --- Discrete Sub-section (Span 8) ---
%     \multicolumn{8}{c}{\textbf{\textit{Discrete Tokenization}}} \\

%     Janus-Pro (7B)        & 2025.01 & 86.9   & 88.9   & 89.4   & 89.3   & 89.5   & 84.2 \\
%     Emu3 (8B)             & 2024.09 & 85.2 & 86.7 & 86.8 & 90.2 & 83.2 & 80.6 \\
%     X-Omni (7B)           & 2025.07 & 84.8   & 92.6   & 90.6   & 94.8   & 84.2   & 87.7 \\
%     Show-O2 (7B)          & 2025.06 &  89.0  & 91.8   & 90.0   & 91.9   & 91.6   & 86.1 \\
%     Unitok (7B) & 2025.02 & --   & --   & --   & --   & --   & -- \\
%     TokenFlow (13B)       & 2024.12 & 78.7 & 79.2 & 81.3 & 85.2 & 71.2 & 73.4 \\
%     Nextflow (7B)         & 2026.01 & 92.4   & 90.1   & 90.5   & 92.7   & 91.1   & 86.0 \\
%     Nextflow-RL (7B)      & 2026.01 & 92.1 & 93.5   & 91.9   & 92.1   & 93.2   & 88.3 \\
%     % \rowcolor{blue!10}
%     %Kelix-TC-DiT (8B)          &       & 85.3 & 91.4 & 89.0 & 73.7 & 79.3 & 84.6 \\
%     \rowcolor{blue!10}
%     % Kelix (8B)          &       & 84.8 & 91.6 & 89.2 & 92.9 & 86.0 & 85.6 \\
%     Kelix (8B)          &       & 84.5 &	91.5 &	89.2 &	93	& 84.4 &	85.5 \\
    
%     \midrule
    
%     \rowcolor{gray!10}
%     \multicolumn{8}{l}{\textbf{Proprietary Models}} \\
%     GPT-4o                & 2024.05 & --   & --   & --   & --   & --   & --   \\

%     \bottomrule
%     \end{tabular}
%     }
%     \label{tab:dpg_bench_results}
% \end{table}
\begin{table}[t]
    \caption{
        Comprehensive comparison of image generation performance on GenEval and WISE benchmarks. The release dates denote the time of arXiv submission or official release. \textbf{For Kelix, we report performance using the default resolution of $1024\times1024$.}
    }
    \centering
    \small
    \renewcommand{\arraystretch}{1.2}
    \resizebox{\textwidth}{!}{
    \begin{tabular}{l c  ccccccc  ccccccc}
    \toprule
    \multirow{2}{*}{\textbf{Model}} & \multirow{2}{*}{\textbf{Date}} & \multicolumn{7}{c}{\textbf{GenEval Benchmark}} & \multicolumn{7}{c}{\textbf{WISE Benchmark}} \\
    \cmidrule(lr){3-9} \cmidrule(lr){10-16}
    & & Single & Two & Count & Color & Pos. & C-Attr & \textbf{Overall} & Cult. & Time & Space & Bio. & Phys. & Chem. & \textbf{Overall} \\
    \midrule
    
    \rowcolor{gray!10}
    \multicolumn{16}{l}{\textbf{Dedicated T2I Models}} \\
    SDXL (3.5B)           & 2023.07   & 98.0 & 74.0 & 39.0 & 85.0 & 15.0 & 23.0 & 55.0 & 43.0 & 48.0 & 47.0 & 44.0 & 45.0 & 27.0 & 43.0 \\
    DALL-E 3              & 2023.10   & 96.0 & 87.0 & 47.0 & 83.0 & 43.0 & 45.0 & 67.0 & --   & --   & --   & --   & --   & --   & --   \\
    SD3 Medium (2B)       & 2024.06   & 99.0 & 94.0 & 72.0 & 89.0 & 33.0 & 60.0 & 74.0 & 43.0 & 50.0 & 52.0 & 41.0 & 53.0 & 33.0 & 45.0 \\
    FLUX.1-dev (12B)      & 2024.08   & 98.0 & 93.0 & 75.0 & 93.0 & 68.0 & 65.0 & 82.0 & 48.0 & 58.0 & 62.0 & 42.0 & 51.0 & 35.0 & 50.0 \\
    \midrule

    \rowcolor{gray!10}
    \multicolumn{16}{l}{\textbf{LLM \& Diffusion Conjunction}} \\
    MetaQuery-XL (7B)     & 2025.05   & --   & --   & --   & --   & --   & --   & 80.0 & 56.0 & 55.0 & 62.0 & 49.0 & 63.0 & 41.0 & 55.0 \\
    OmniGen-2 (7B)        & 2025.06   & 100.0 & 95.0 & 64.0 & 88.0 & 55.0 & 76.0 & 80.0 & --   & --   & --   & --   & --   & --   & --   \\
    Qwen-Image (27B)      & 2025.08   & 99.0 & 92.0 & 89.0 & 88.0 & 76.0 & 77.0 & 87.0 & -- & -- & -- & -- & -- & -- & -- \\
    \midrule

    \rowcolor{gray!10}
    \multicolumn{16}{l}{\textbf{Unified Multimodal LLMs}} \\
    
    % --- Continuous Sub-section (Centered) ---
    \multicolumn{16}{c}{\textbf{\textit{Continuous Tokenization}}} \\
    Bagel (14B)           & 2025.05   & 99.0 & 94.0 & 81.0 & 88.0 & 64.0 & 63.0 & 82.0 & 44.0 & 55.0 & 68.0 & 44.0 & 60.0 & 39.0 & 52.0 \\
    Ming-Univision (16B)  & 2025.10   & 100.0   & 93.0   & 59.0   & 93.0   & 92.0   & 70.0   & 85.0 & --   & --   & --   & --   & --   & --   & --   \\
    Blip-3o (4B)          & 2025.05   & --   & --   & --   & --   & --   & --   & 81.0 & --   & --   & --   & --   & --   & --   & 50.0 \\
    TUNA (7B)             & 2025.12   & 100   & 97   & 81   & 91   & 88   & 83   & 90.0 & --   & --   & --   & --   & --   & --   & --   \\

    % --- Hybrid Sub-section (Centered) ---
    \multicolumn{16}{c}{\textbf{\textit{Hybrid Tokenization}}} \\
    Manzano (30B)         & 2025.09   & 100.0 & 91.0 & 83.0 & 87.0 & 84.0 & 65.0 & 85.0 & 58.0 & 50.0 & 65.0 & 50.0 & 55.0 & 32.0 & 54.0 \\
    VQ-RAE (7B)  & 2025.11   & 96.0 & 82.0 & 64.0 & 80.0 & 73.0 & 58.0 & 76.0 & -- & -- & -- & -- & -- & -- & -- \\

    % --- Discrete Sub-section (Centered) ---
    \rowcolor{blue!10}
    \multicolumn{16}{c}{\textbf{\textit{Discrete Tokenization}}} \\
    \rowcolor{blue!10}
    Janus-Pro (7B)        & 2025.01   & \underline{99.0} & 89.0 & 59.0 & 90.0 & \underline{79.0} & 66.0 & 80.0 & 30.0 & 37.0 & 49.0 & 36.0 & 42.0 & 26.0 & 35.0 \\
    \rowcolor{blue!10}
    Emu3 (8B)             & 2024.09   & \underline{99.0}   & 81.0   & 42.0   & 80.0   & 49.0   & 45.0   & 66.0 & --   & --   & --   & --   & --   & --   & 39.0 \\
    \rowcolor{blue!10}
    X-Omni (7B)           & 2025.07   & 98.0 & \textbf{95.0} & \underline{75.0} & \underline{91.0} & 71.0 & 68.0 & 83.0 & --   & --   & --   & --   & --   & --   & --   \\
    \rowcolor{blue!10}
    Show-O2 (7B)          & 2025.06   & \textbf{100.0}   & 87.0   & 58.0   & \textbf{92.0}   & 52.0   & 62.0 & 76.0 & --   & --   & --   & --   & --   & --   & --   \\
    \rowcolor{blue!10}
    Unitok (7B)           & 2025.02   & \underline{99.0}   & 71.0   & 36.0   & 79.0   & 26.0   & 45.0   & 59.0 & --   & --   & --   & --   & --   & --   & --   \\
    \rowcolor{blue!10}
    TokenFlow (13B)       & 2024.12   & 97.0   & 66.0   & 40.0   & 84.0   & 17.0   & 26.0   & 55.0 & --   & --   & --   & --   & --   & --   & --   \\
    \rowcolor{blue!10}
    Nextflow (7B)         & 2026.01   & 98.0   & 92.0   & 73.0   & 90.0   & 77.0   & \underline{69.0}   & 83.0 & \textbf{62.0}   & \underline{60.0}   & \textbf{70.0}   & \textbf{54.0}   & \underline{58.0}   & \textbf{38.0}   & \textbf{59.0} \\
    \rowcolor{blue!10}
    % Nextflow-RL (7B)      & 2026.01   & \textbf{100.0}  & 92.0   & \underline{75.0}   & 90.0   & 76.0   & \underline{70.0}   & \underline{84.0} & \textbf{63.0}   & \textbf{63.0}   & \textbf{77.0}   & \textbf{58.0}   & 67.0   & \textbf{39.0}   & \textbf{62.0} \\
    \rowcolor{blue!10}
    Kelix (8B) & & \textbf{100.0} & \underline{93.0} & \textbf{81.0} & 87.0 & \textbf{85.0} & \textbf{79.0} & \textbf{87.6} & \underline{58.0} & \textbf{62.0} & \underline{68.0} & \underline{47.0} & \textbf{61.0} & \underline{35.0} & \underline{57.0} \\
    \midrule
    
    \rowcolor{gray!10}
    \multicolumn{16}{l}{\textbf{Proprietary Models}} \\
    GPT-4o                & 2024.05   & 99.0 & 92.0 & 85.0 & 92.0 & 75.0 & 61.0 & 84.0 & 81.0 & 71.0 & 89.0 & 83.0 & 79.0 & 74.0 & 80.0 \\
    
    \bottomrule
    \end{tabular}
    }
    \label{tab:generation_results}
\end{table}


% \begin{table}[t!]
%     \caption{
%         Comparison of image generation performance on the DPG-Bench benchmark. All scores are reported on a 100-point scale. The release dates denote the time of arXiv submission or official release.
%     }
%     \centering
%     \small
%     \renewcommand{\arraystretch}{1.2}
%     \resizebox{0.9\textwidth}{!}{
%     \begin{tabular}{l c  ccccc  c}
%     \toprule
%     \textbf{Method} & \textbf{Date} & \textbf{Global} & \textbf{Entity} & \textbf{Attribute} & \textbf{Relation} & \textbf{Other} & \textbf{Overall$\uparrow$} \\
%     \midrule
    
%     \rowcolor{gray!10}
%     \multicolumn{8}{l}{\textbf{Dedicated T2I Models}} \\
%     SDXL (2.6B)           & 2023.07 & 83.3 & 82.4 & 80.9 & 86.8 & 80.4 & 74.7 \\
%     DALL-E 3              & 2023.10 & 91.0 & 89.6 & 88.4 & 90.6 & 89.8 & 83.5 \\
%     SD3 Medium (2B)       & 2024.06 & 87.9 & 91.0 & 88.8 & 80.7 & 88.7 & 84.1 \\
%     FLUX.1-dev (12B)      & 2024.08 & --   & --   & --   & --   & --   & --   \\
%     \midrule

%     \rowcolor{gray!10}
%     \multicolumn{8}{l}{\textbf{LLM \& Diffusion Conjunction}} \\
%     MetaQuery-XL (7B)     & 2025.04 & --   & --   & --   & --   & --   & 82.05  \\
%     OmniGen-2 (7B)        & 2025.06 & 88.81   & 88.83   & 90.18   & 89.37   & 90.27   & 83.57   \\
%     Qwen-Image (27B)      & 2025.08 & 91.32   & 91.56   & 92.02   & 94.31   & 92.73   & 88.32   \\
%     \midrule

%     \rowcolor{gray!10}
%     \multicolumn{8}{l}{\textbf{Unified Multimodal LLMs}} \\
    
%     \multicolumn{8}{c}{\textbf{\textit{Continuous Tokenization}}} \\
%     Bagel (14B)           & 2025.05 & --   & --   & --   & --   & --   & --   \\
%     Ming-Univision (16B)  & 2025.10 & --   & --   & --   & --   & --   & 82.1 \\
%     Blip-3o (4B)          & 2025.05 & --   & --   & --   & --   & --   & 79.4 \\
%     TUNA (7B)             & 2025.12 & 90.42   & 91.58   & 90.94   & 91.87   & 90.73   & 86.8 \\

%     \multicolumn{8}{c}{\textbf{\textit{Hybrid Tokenization}}} \\
%     Manzano (30B)         & 2025.09 & --   & --   & --   & --   & --   & --   \\

%     \multicolumn{8}{c}{\textbf{\textit{Discrete Tokenization}}} \\
%     Janus-Pro (7B)        & 2025.01 & 86.9   & 88.9   & 89.4   & 89.3   & 89.5   & 84.2 \\
%     Emu3 (8B)             & 2024.09 & 85.2 & 86.7 & 86.8 & 90.2 & 83.2 & 80.6 \\
%     X-Omni (7B)           & 2025.07 & 84.8   & 92.6   & 90.6   & 94.8   & 84.2   & 87.7 \\
%     Show-O2 (7B)          & 2025.06 &  89.0  & 91.8   & 90.0   & 91.9   & 91.6   & 86.1 \\
%     Unitok (7B)           & 2025.02 & --   & --   & --   & --   & --   & -- \\
%     TokenFlow (13B)       & 2024.12 & 78.7 & 79.2 & 81.3 & 85.2 & 71.2 & 73.4 \\
%     Nextflow (7B)         & 2026.01 & 92.4   & 90.1   & 90.5   & 92.7   & 91.1   & 86.0 \\
%     Nextflow-RL (7B)      & 2026.01 & 92.1 & 93.5   & 91.9   & 92.1   & 93.2   & 88.3 \\
%     \rowcolor{blue!10}
%     Kelix (8B)          &       & 84.5 & 91.5 & 89.2 & 93 & 84.4 & 85.5 \\
    
%     \midrule
    
%     \rowcolor{gray!10}
%     \multicolumn{8}{l}{\textbf{Proprietary Models}} \\
%     GPT-4o                & 2024.05 & --   & --   & --   & --   & --   & --   \\

%     \bottomrule
%     \end{tabular}
%     }
%     \label{tab:dpg_bench_results}
% \end{table}

\begin{table}[t!]
\footnotesize
    \caption{
        Comparison of image generation performance on the DPG-Bench benchmark. All scores are reported on a 100-point scale. The release dates denote the time of arXiv submission or official release.
    }
    \centering
    \small
    \renewcommand{\arraystretch}{1.2}
    % 如果表格略微超过页宽，可以取消下面这行的注释来微调列间距
    % \setlength{\tabcolsep}{5pt} 
    
    \begin{tabular}{l c  ccccc  c}
    \toprule
    \textbf{Method} & \textbf{Date} & \textbf{Global} & \textbf{Entity} & \textbf{Attribute} & \textbf{Relation} & \textbf{Other} & \textbf{Overall$\uparrow$} \\
    \midrule
    
    \rowcolor{gray!10}
    \multicolumn{8}{l}{\textbf{Dedicated T2I Models}} \\
    SDXL (2.6B)           & 2023.07 & 83.3 & 82.4 & 80.9 & 86.8 & 80.4 & 74.7 \\
    DALL-E 3              & 2023.10 & 91.0 & 89.6 & 88.4 & 90.6 & 89.8 & 83.5 \\
    SD3 Medium (2B)       & 2024.06 & 87.9 & 91.0 & 88.8 & 80.7 & 88.7 & 84.1 \\
    FLUX.1-dev (12B)      & 2024.08 & --   & --   & --   & --   & --   & --   \\
    \midrule

    \rowcolor{gray!10}
    \multicolumn{8}{l}{\textbf{LLM \& Diffusion Conjunction}} \\
    MetaQuery-XL (7B)     & 2025.04 & --   & --   & --   & --   & --   & 82.05  \\
    OmniGen-2 (7B)        & 2025.06 & 88.81   & 88.83   & 90.18   & 89.37   & 90.27   & 83.57   \\
    Qwen-Image (27B)      & 2025.08 & 91.32   & 91.56   & 92.02   & 94.31   & 92.73   & 88.32   \\
    \midrule

    \rowcolor{gray!10}
    \multicolumn{8}{l}{\textbf{Unified Multimodal LLMs}} \\
    
    % --- Continuous Sub-section (Span 8) ---
    \multicolumn{8}{c}{\textbf{\textit{Continuous Tokenization}}} \\
    Bagel (14B)           & 2025.05 & --   & --   & --   & --   & --   & --   \\
    Ming-Univision (16B)  & 2025.10 & --   & --   & --   & --   & --   & 82.1 \\
    Blip-3o (4B)          & 2025.05 & --   & --   & --   & --   & --   & 79.4 \\
    TUNA (7B)             & 2025.12 & 90.42   & 91.58   & 90.94   & 91.87   & 90.73   & 86.8 \\

    % --- Hybrid Sub-section (Span 8) ---
    \multicolumn{8}{c}{\textbf{\textit{Hybrid Tokenization}}} \\
    Manzano (30B)         & 2025.09 & --   & --   & --   & --   & --   & --   \\
    VQ-RAE (7B)         & 2025.11 & 89.8   & 93.1   & 89.9   & 90.3   & 91.2   & 86.6   \\

    % --- Discrete Sub-section (Span 8) ---
    \rowcolor{blue!10}
    \multicolumn{8}{c}{\textbf{\textit{Discrete Tokenization}}} \\
    \rowcolor{blue!10}
    Janus-Pro (7B)        & 2025.01 & 86.9   & 88.9   & 89.4   & 89.3   & 89.5   & 84.2 \\
    \rowcolor{blue!10}
    Emu3 (8B)             & 2024.09 & 85.2 & 86.7 & 86.8 & 90.2 & 83.2 & 80.6 \\
    \rowcolor{blue!10}
    X-Omni (7B)           & 2025.07 & 84.8   & \textbf{92.6}   & \textbf{90.6}   & \textbf{94.8}   & 84.2   & \textbf{87.7} \\
    \rowcolor{blue!10}
    Show-O2 (7B)          & 2025.06 &  \underline{89.0}  & \underline{91.8}   & 90.0   & 91.9   & \textbf{91.6}   & \underline{86.1} \\
    \rowcolor{blue!10}
    Unitok (7B)           & 2025.02 & --   & --   & --   & --   & --   & -- \\
    \rowcolor{blue!10}
    TokenFlow (13B)       & 2024.12 & 78.7 & 79.2 & 81.3 & 85.2 & 71.2 & 73.4 \\
    \rowcolor{blue!10}
    Nextflow (7B)         & 2026.01 & \textbf{92.4}   & 90.1   & \underline{90.5}   & 92.7   & \underline{91.1}   & 86.0 \\
    \rowcolor{blue!10}
    % Nextflow-RL (7B)      & 2026.01 & 92.1 & 93.5   & 91.9   & 92.1   & 93.2   & 88.3 \\
    \rowcolor{blue!10}
    Kelix (8B)          &       & 84.5 & 91.5 & 89.2 & \underline{93.0} & 84.4 & 85.5 \\
    
    % \midrule
    
    % \rowcolor{gray!10}
    % \multicolumn{8}{l}{\textbf{Proprietary Models}} \\
    % GPT-4o                & 2024.05 & --   & --   & --   & --   & --   & --   \\

    \bottomrule
    \end{tabular}
    \label{tab:dpg_bench_results}
\end{table}

\section{Training}

In this section, we describe the dataset construction, stage-wise training configurations, evaluation benchmark results, and detailed ablation studies.



\subsection{Datasets}
We curate a large-scale, diverse, and high-quality corpus exceeding 1T tokens to train the full Kelix pipeline, spanning the discrete image tokenizer, the unified LLM, and the diffusion-based image de-tokenizer. Our training data covers the following categories:

%
\begin{itemize}
    \item \textbf{Image Caption}: Large-scale image-text pairs that establish foundational world knowledge by mapping visual features to linguistic concepts. Sources include Coyo~\citep{kakaobrain2022coyo-700m}, DataComp~\citep{gadre2023datacomp}, LAION~\citep{schuhmann2022laion}, Infinity-MM Stage1~\citep{gu2024infinity}, CC12M~\citep{changpinyo2021conceptual}, PD12M~\citep{meyer2024public}, ConceptualCaptions~\citep{sharma2018conceptual}, DenseFusion~\citep{li2024DenseFusion}, among others.
    \item \textbf{Optical Character Recognition}: OCR data spanning handwritten text, charts, synthetic documents, and multi-image layouts, sourced from OneVision~\citep{zheng2025onevision}, CodeOCR~\citep{codeocr_leetcode_2025}, SynthDoG~\citep{kim2022ocr}, InstructOCR~\citep{duan2025instructocr}, CCPD~\citep{xu2018towards}, and others. To address the scarcity of Chinese OCR data, we synthesize in-house datasets with multi-turn QA tasks covering diverse character sets, non-conventional reading orders, and specialized layouts (\textit{e.g.}, ancient Chinese poems), thereby strengthening the model's ability to parse spatial arrangement, font variation, and instruction-following.To enhance the model’s OCR grounding and performance on text-dense tasks, we apply rejection sampling on natural images to select samples that are (i) text-heavy, (ii) contain small and hard-to-read text, and (iii) exhibit sparsely and discretely distributed text regions. Using Qwen3-VL and Gemini, we then construct multi-dimensional OCR-focused QA pairs. In addition, since document OCR capability is critical, we incorporate training data spanning multiple programming languages as well as HTML, PDF, and Markdown formats.

\begin{figure}[t!]
  \centering
  \includegraphics[width=0.9\textwidth]{OneBase/figures/generation_data_pipeline.pdf}
    \caption{Overview of the high-quality text-to-image training data synthesis pipeline, comprising prompt construction, frontier model distillation, VLM-based multi-dimensional quality verification, and final selection.}
    \label{fig:generation_data_pipeline}
\end{figure}
    
    \item \textbf{Visual Question Answering}: VQA data from diverse domains to strengthen context-aware visual reasoning, including ShareGPT-4o~\citep{cui2025comprehensive}, Docmatix~\citep{laurencon2024building}, M4-instruct~\citep{li2024llavanext-interleave}, ShareGPT4V~\citep{chen2024sharegpt4v}, DVQA~\citep{kafle2018dvqa}, Mavis-Function~\citep{zhang2024mavis}, Honey-Bee~\citep{bansal2025honeybee}, FineVision~\citep{wiedmann2025finevision}, etc.
    \item \textbf{STEM \& Interleaved}: Interleaved image-text data for long-context multimodal modeling, drawn from Wikipedia~\citep{wikidump}, ArxivPapers, MM-Interleaved~\citep{tian2024mm}, and Pdf2Doc.
    \item \textbf{Pure Text}: A large-scale text corpus spanning math, code, chain-of-thought reasoning, and common knowledge, used to preserve the backbone's pretrained capabilities. Sources include Nemotron-CC~\citep{nvidia_nemotron_nano_v3_2025}, Nemotron-CC-Math~\citep{karimi2025nemotroncc,nvidia2025nvidianemotronnano2}, GeneralThought-430K~\citep{RJT1990_2025GeneralThoughtArchive}, NuminaMath-QwQ-CoT-5M~\citep{primeintellect2024numinamath}, and OpenCodeReasoning~\citep{ahmad2025opencodereasoning}.
    % \item \textbf{Text-to-Image}: For image generation, we collect text-to-image paired data from CommonCatalog~\citep{gokaslan2024commoncanvas}, midjourney-niji-1m-llavanext~\citep{midjourney-niji-1m-llavanext}, BLIP3o-Pretrain-Short-Caption and BLIP3o-Pretrain-Long-Caption~\citep{chen2025blip3}, along with inverted caption datasets. Beyond public sources, we synthesize in-house data targeting fine-grained control over object attributes, spatial relations, counting, and text rendering. Specifically, we construct seed prompts at multiple complexity tiers (rule-based short prompts and LLM-generated medium/long prompts), generate candidate images via frontier open-source models (Qwen-Image) using best-of-N sampling, and apply multi-dimensional VLM-based quality verification (text--image alignment, attribute accuracy, text fidelity, style, and safety) using Qwen2.5-VL-72B. Rejected images with acceptable visual quality are recaptioned to recover valid training pairs. An overview of the data construction pipeline is presented in Figure \ref{fig:generation_data_pipeline}.
    \item \textbf{Text-to-Image}: For image generation, we curate text-to-image paired data from CommonCatalog~\citep{gokaslan2024commoncanvas}, midjourney-niji-1m-llavanext~\citep{midjourney-niji-1m-llavanext}, BLIP3o-Pretrain-Short-Caption and BLIP3o-Pretrain-Long-Caption~\citep{chen2025blip3}, along with inverted caption datasets. Beyond public sources, we synthesize in-house data targeting fine-grained control over object attributes, spatial relations, counting, and text rendering. As illustrated in Figure~\ref{fig:generation_data_pipeline}, we first construct seed prompts at multiple complexity tiers (rule-based short prompts and LLM-generated medium/long prompts), then generate candidate images via frontier open-source models (Qwen-Image) with best-of-N sampling, and finally apply multi-dimensional VLM-based quality verification (text--image alignment, attribute accuracy, text fidelity, style, and safety) using Qwen2.5-VL-72B. Rejected images with acceptable visual quality are recaptioned to recover valid training pairs.
\end{itemize}


\subsection{Kelix-Tokenizer Training Recipes}

\begin{figure}[t!]
  \centering
  \includegraphics[width=0.75\textwidth]{OneBase/figures/token-recipes.png}
    \caption{Semantic tokenizer training recipe. In Stage 1, we freeze the ViT and LLM backbones. To ensure that the commitment loss can be optimized, we insert the MLP and SimVQ up-projector before VQ codebook. In Stage2 and Stage3, we enable the gradients of ViT and Qwen3.}
    \label{fig:tok-recipe}
\end{figure}


Figure \ref{fig:tok-recipe} shows the training recipes of Kelix-Tok.
The backbone of Kelix-Tok is initialized from Keye-VL 1.5~\citep{team2025kwai}'s pre-trained naitive resolution ViT (NaViT). To enhance the vision-language semantic alignment, we connect the ViT to a small LLM, i.e., Qwen3-0.6B ~\citep{yang2025qwen3}.
% 0.6B Qwen3 model through a single multiple discrete adapter and train it via next-token prediction. 
% To ensure stable codebook training, we employ the SimVQ training strategy.  We introduce a meticulously designed three-stage training process to optimize the tokenizer effectively. 
%
%For the codebook initialization, we first collect a large-scale image corpus to produce patch embeddings by NaVIT, and then conduct the Kmeans algorithm on the patch embeddings to obtain the centroid embedding as the VQ codebook representations.
%
%Next, for stable and fast convergency, we freeze the Kmeans centroid embeddings and apply the SimVQ~\citep{zhu2025addressing} technique to train the codebook.
%
\subsubsection{Codebook Initialization}
For VQ codebook initialization, we first leverage the pre-trained NaViT to extract patch embeddings from a large-scale image corpus, then perform K-means clustering on these embeddings to obtain high-quality cluster centers as the initial representations of the VQ codebook. These cluster centers are further randomly partitioned into N equal size non-overlapping sub-codebooks to ensure complementary semantic coverage across different subspaces. To achieve stable and fast convergence of codebook training, we freeze the centroid embeddings derived from K-means clustering and adopt the SimVQ~\citep{zhu2025addressing} training strategy to optimize the codebook. This strategy only updates the per-subspace projection matrices, which maximizes codebook utilization and effectively mitigates the codebook collapse problem.

\subsubsection{Training Pipline}
Kelix-Tok is optimized through a meticulously designed three-stage training pipeline, with each stage tailored to distinct training objectives to progressively boost the tokenizer's cross-modal semantic encoding capability, as detailed below:

\paragraph{Stage 1: Codebook Alignment} % This stage isolates the multiple discrete adapters for training in image captions, while the core NaViT and LLM components remain frozen. This focused phase allows the codebook to establish an initial visual-semantic mapping. % during which we observe a compelling ``epiphany'' phenomenon. 
This stage solely trains the multi-discrete adapters with image caption corpora, while the core NaViT encoder and Qwen3-0.6B LLM backbone remain frozen. This stage enables the VQ codebook to establish an initial visual-semantic mapping, laying a fundamental foundation for subsequent cross-modal alignment between visual and linguistic representations.

\paragraph{Stage 2: Multi-Task Pre-training} % In this stage, we expand the scope to full-parameter training across a rich dataset of understanding tasks (general knowledge, reasoning, text-rich) to enhance the tokenizer's comprehensive abilities. Our objective is to expose the tokenizer to the broadest possible data distribution to explore the storage capacity of the codebook and expand its knowledge reservoir; meanwhile, we also adapt the tokenizer to better autoregressive generation paradigm of the large language model (LLM) at this stage.
We unlock all model parameters for end-to-end training on a diverse collection of understanding-oriented datasets, covering general knowledge, logical reasoning, and text-rich tasks. By exposing the tokenizer to a broad data distribution, this stage fully explores the storage capacity of the codebook and expands its knowledge reservoir. Meanwhile, it adapts the tokenizer to the autoregressive generation paradigm of LLMs, achieving deep alignment between discrete visual representations and the LLM's semantic space.

\paragraph{Stage 3: Annealing Training} % The last stage fine-tunes the tokenizer on high-quality instruction-tuning datasets. % with a low learning rate. 
% At this stage, our objective is to enable the tokenizer to capture more inherent rules and latent information from the data.
As the final fine-tuning stage, we optimize the tokenizer on high-quality instruction-tuning datasets with an annealing learning rate. The core objective is to capture the inherent underlying rules and latent semantic information from the data, further enhancing the expressiveness of discrete visual representations and the robustness of cross-modal reasoning.

Upon the completion of the entire training pipeline, the auxiliary Qwen3-0.6B LLM decoder is discarded. The resulting Kelix-Tok, an integration of the NaViT encoder and a set of VQ codebooks, serves as a dedicated image tokenizer that encodes raw images into structured discrete visual representations. This design allows the knowledge of the image domain to be seamlessly transferred to various LLM decoders, realizing efficient cross-modal knowledge sharing. % for downstream unified multimodal modeling.

% Once training is complete, the LLM decoder is discarded. The Kelix-Tok, which integrates a Vision Transformer (ViT) and a set of vector quantization (VQ) codebooks, then serves as a dedicated tokenizer to encode the image modality. 
% It converts the image into discrete representation, so that knowledge of the image domain can be seamlessly transferred to other LLM decoders.
%, converting the image modality into discrete input and output representation for LLM, which can be seamlessly transferred to other LLM decoders with minnor training adaptation.


\subsection{Kelix-LLM Training Recipes}
%\renewcommand{\leftmargin}{\parindent}
\setlength{\leftmargin}{\parindent}
\setlength{\itemindent}{0pt}
\setlength{\labelwidth}{0pt}
\label{sec:training_strategy}

\begin{figure}[t!]
  \centering
  \includegraphics[width=0.75\textwidth]{OneBase/figures/AR-recipes-0210-2.png}
    \caption{The LLM training recipe of Kelix, we first pre-train the \textit{Block Decoder} and the image token embedding for a better warm start, and then unifying the image token and text token to a same embedding table for full-parameter tuning.}
    \label{fig:ar-recipe}
\end{figure}



The proposed unified model follows a four-stage progressive training strategy to bridge textual and visual modalities while preserving the core reasoning capabilities of the LLM. 
This LLM training protocol progressively optimizes different model components, with data scale and composition  tailored to the objectives of each stage. 
%It starts from cross-modal alignment using moderate-scale caption data, then expands to large-scale interleaved image-text and multi-task pre-training corpora. 
%The model is subsequently refined via annealing-based training with Chain-of-Thought (CoT) data and high-quality filtered datasets, and finally undergoes task-specific supervised fine-tuning (SFT) using specialized CoT data and targeted image generation samples. 
%Besides, pure text corpus are also used across different training stages to preserve Qwen3-pretrained linguistic abilities.
Table \ref{tab:ar_training_data} summarizes the number of samples and token volume for each training stage, illustrating the progressive increase in data complexity and the shift toward high-quality, task-specific data in later stages.

\subsubsection{Visual Processor}
To enhance spatial position awareness under Qwen3’s 1D Rotary Position Embedding (RoPE) framework, we introduce a dedicated visual processor based on Keye-VL1.5~\citep{yang2025kwai}, uniformly applied to both image understanding and generation.

%The visual processor’s core mechanism: After the native-resolution NaViT encoder compresses images into discrete patch tokens, each token row is prepended with task-agnostic special tokens and embedding "row, column" spatial coordinates (e.g., <|mm\_pos\_start|>3,4<|mm\_pos\_end|> for the 3rd row of a 4-column grid). This supplements 1D RoPE with explicit 2D spatial cues, enabling effective inter-patch positional modeling without modifying the transformer core.
With the patch tokens given by Kelix-Tok, each token row is prepended with special tokens and embedded with ``row, column'' spatial coordinates (\textit{e.g.}, <|mm\_pos\_start|>3,4<|mm\_pos\_end|>). These markers extend 1D RoPE with structured 2D spatial information, allowing accurate modeling of inter-patch positional relationships while preserving the original Transformer architecture.

%For image generation, the diffusion-based de-tokenizer (DiT) accepts a fixed 442-length conditional embedding from the AR model’s hidden states, constructed as:
%- 324 image tokens: Encoded from 504×504 images (constrained in training) via an 18×18 patch grid (504/28=18) for uniform spatial sampling;
%- 118 position encoding tokens: Reinforce global positional awareness for DiT’s spatial reasoning;
%- Concatenation of the above two components forms the 442-length embedding, fed into DiT’s y-embedder to guide semantic-aligned generation, ensuring consistent AR-DiT alignment.



\subsubsection{Stage 1: Cross-Modal Alignment}
This stage serves as a ``warm-up'' to establish a stable mapping between discrete visual representations and the LLM’s linguistic latent space, avoiding catastrophic forgetting of pre-trained linguistic priors. All images are processed at native resolution, as illustrated in Stage 1 on the left of Figure \ref{fig:ar-recipe}.

In \textit{Block Encoder}, the textual modality is mapped to textual embeddings via \textit{Qwen3-8B Embedding}; the visual modality is encoded into VQ embeddings by Kelix-Tok. To match the LLM backbone’s dimensionality, the VQ embedding is transformed via a linear \textit{up-projector}; Decoder layers and embedding layer of Qwen3-8B are fully frozen; only the \textit{Block Decoder} and LM head are optimized, with training data comprising large-scale pure text corpora and image-caption pairs. % aligned with the unified description style.

This stage lays the foundation for cross-modal semantic consistency by completing the initial alignment between visual and textual representations.

\subsubsection{Stage 2: Unified Multimodal Pre-Training}
Building on Stage 1, this stage integrates embedding space unification, cross-modal generalization scaling, and multi-task capability development. It transforms the LLM into a native multimodal engine with integrated understanding, reasoning, and generation capabilities through a two-substage progressive training process.

\paragraph{Embedding Unification}
Following Stage 1, the VQ embedding table is projected sequentially by two pre-trained \textit{up-projectors} to generate a unified image embedding table: the first is the pre-trained \textit{up-projector} from the Kelix-Tok, and the second is the \textit{up-projector} fine-tuned in LLM Stage 1. This image embedding table is then concatenated with Qwen3’s text embedding table to form a unified \textit{Block Encoder}, which replaces the separate Qwen3 Embedding, VQ Embedding, and \textit{up-projector} components used in Stage 1. 

\paragraph{Interleaved Scaling (Stage 2.1)}
This substage focuses on achieving seamless integration of textual and visual modalities by unifying the embedding space, while expanding data diversity to enhance cross-modal generalization. All model parameters are unfrozen to enable end-to-end autoregressive training.
%
% \textbf{Block Embedding Unification} 
Training data utilizes large-scale interleaved image-text data and text corpora to balance cross-modal scaling and linguistic capability preservation.

This substage lays the structural and data foundation for complex multi-task learning by unifying the embedding space and expanding cross-modal data coverage.

\paragraph{Multi-Task Pre-Training (Stage 2.2) }
Building on the previous substage, this substage enriches the model’s multimodal capabilities through diverse task data, shaping basic understanding, reasoning, and generation competencies. Training relies on a large-scale, diverse corpus covering pure text corpora, descriptive captions, simple visual-language QA, and text-to-image pairs.
%
For image generation tasks, input images are constrained to an aspect ratio range of 0.75 to 1.33 and then resized to a fixed resolution of 504×504 pixels. We also enable the completion-only loss mask that backpropagates negative log-likelihood loss solely through the assistant’s response tokens. % excluding prompt tokens.

This substage equips the model with comprehensive multimodal capabilities by integrating diverse task data, balancing understanding and generation performance for subsequent refinement.

\subsubsection{Stage 3: Annealing Training}
This stage focuses on enhancing instruction following, logical reasoning, and generative quality, with learning rate annealing initiated to ensure stable convergence. Curated high-quality training data includes CoT data, high-quality filtered multimodal QA, filtered text corpora, and general image generation data, with rejection sampling applied to filter low-quality data, prevent hallucinations, and enhance fine-grained abilities.

This stage refines the model’s reasoning rigor and generative stability, laying the foundation for task-specific optimization.

\subsubsection{Stage 4: Supervised Fine-Tuning (SFT)}
Building directly on Stage 3, this stage retains the annealed learning rate and focuses on optimizing task-specific performance while preserving unified multimodal capabilities. Core training data includes replayed high-quality Stage 3 data, task-specific CoT data, targeted image generation data, and task-specific enhancement data, with rejection sampling retained to filter noisy samples and further mitigate hallucinations.

This stage boosts the model’s performance on specific tasks while maintaining the balance between understanding and generation, ensuring practical utility in real-world scenarios.

\begin{table}[t!]
\centering
\footnotesize % 保持小字体适配紧凑排版
\renewcommand{\arraystretch}{1.4} % 优化行间距
\setlength{\tabcolsep}{2pt} % 维持最小列间距
\begin{tabular}{p{4cm}ccp{9cm}} % 适配数据类型列换行
\toprule
Stage & Samples & Tokens & Data Type \\
\midrule
Stage 1:\newline Cross-Modal Alignment & 191M & 108B & Pure text corpora; \newline image-caption pairs  % (\textit{e.g.}, Coyo, LAION, CC12M) 
\\
\midrule
Stage 2.1:\newline Interleaved Scaling & 173M & 191B & % Large-scale interleaved \newline image-text data \\
Large-scale interleaved image-text data \\
% \midrule
Stage 2.2:\newline Multi-Task Pre-Training & 593M & 679B & Large-scale pure text corpora; \newline descriptive captions; simple visual-language QA; \newline text-to-image pairs (reversed from captions) \\
\midrule
Stage 3:\newline Annealing Training & 40M & 70B & Pure text Chain-of-Thought (CoT) data; \newline multimodal CoT data; high-quality filtered multimodal QA; \newline filtered pure text corpora; general image generation data (filtered via rejection sampling) \\
\midrule
Stage 4:\newline SFT & 4M & 8B & Replayed Stage 3 high-quality data; \newline task-specific CoT data; task-specific image generation data (e.g., attribute/subject/scene/style/perspective control); \newline task-specific enhancement data (\textit{e.g.}, complex OCR, math reasoning, STEM data) \\
\bottomrule
\end{tabular}
\caption{Data statistics of each LLM training stage, including the number of samples, total token volume, and data type.}
\label{tab:ar_training_data}
\end{table}


\subsection{Kelix-DiT Training Recipes}
We exclusively target 1024×1024 high-resolution image generation and adopt a two-stage training paradigm, which balances the model’s generalization ability across diverse scenarios and fine-grained adherence to textual instructions.



\subsubsection{Training Setting}
Kelix-DiT is conditioned on all last hidden states between the vision start token and the vision end token.
% Kelix-DiT is conditioned on a fixed 442-token embedding derived from the LLM's hidden representations. This condition consists of (1) 324 image tokens encoded from 504$\times$504 inputs using an 18$\times$18 patch grid, and (2) 118 positional encoding tokens that provide global spatial context. % The combined embedding is passed to DiT's y-embedder, enabling semantically coherent generation and stable alignment between AR and diffusion components. 
% For conditional feature extraction, input images are resized to 504×504 to match the input requirement of the pre-trained unified LLM model. Consequently, the DiT receives conditional signals encoded from the AR model’s hidden states, with a fixed sequence length of 442. For the target image, it is resized to 1024×1024 and encoded into 32×32 latent representations via a pre-trained DC-AE-F32C32\citep{he2025dc} (abbreviated as DC-AE), which employs a 32× downsampling ratio to compress pixel-level information into a compact latent space. 
All DiT parameters are updated during training.

% During the entire training process, only the parameters of the DiT backbone and the dedicated connector module (for bridging AR hidden states and DiT conditional inputs) are updated. The LLM, DC-AE encoder, and DC-AE decoder are kept frozen to preserve their pre-trained semantic encoding capabilities. % and latent manifold modeling properties, avoiding interference from the detokenizer training process.

\subsubsection{Two-Stage Training Strategy}
\paragraph{Pretraining Stage} The pretraining dataset is constructed by inverting large-scale image-caption pairs, with source images constrained to an aspect ratio of 0.67–1.5. These images are then resized to 504×504 for condition extraction and 1024×1024 for target construction. The dataset’s diversity in scenes, object categories, and image scales equips the model with robust semantic-image alignment capabilities and strong generalization to unseen scenarios.


\paragraph{SFT Stage} We use high-quality text-to-image data filtered via rejection sampling to enhance the model’s instruction-following capability. The dataset is carefully curated to cover key fine-grained control dimensions, including attribute control, subject control and so on. The target image is strictly constrained to a resolution of 1024x1024.
% including attribute control (e.g., color, size), subject control (e.g., quantity, type), scene generation (e.g., indoor/outdoor, time of day), artistic style control (e.g., oil painting, sketch), and perspective control (e.g., top-down view, close-up). This targeted training facilitates the generation of high-fidelity images that strictly comply with diverse textual instructions.


% \subsection{Kelix-Tok-DiT Training Recipe}
% % 
% The training pipeline of Kelix-Tok-DiT consists of three core components: (1) the pre-trained Kelix-Tok (detailed in \ref{sec:discrete_image_tokenizer}) for discrete visual token extraction; (2) a diffusion transformer (DiT) backbone (i.e., SANA 1.5\citep{xie2024sana}) for latent denoising and image synthesis; and (3) the DC-AE-F32C32 model (abbreviated as DC-AE\citep{he2025dc}) for latent space mapping.

% We enable native-resolution encoding for condition images, meaning no cropping or resizing operations are applied to input images. This design preserves the complete spatial structure and semantic information of the original visual content, which is critical for capturing fine-grained details in generated images. For target image construction, original images are center-cropped and resized to a fixed resolution of 512×512 to ensure consistent training objectives. During this stage, only the DiT backbone is optimized, while the Kelix-Tok and DC-AE components remain frozen. 
% to retain their pre-trained capabilities.




\section{Evaluation}

\subsection{Benchmarks}
In this section, we evaluate our model’s image understanding and generation capabilities across diverse benchmarks, covering core task scenarios and ensuring rigorous assessment.

\subsubsection{Image Understanding Benchmarks}
We categorize the evaluation into three types of benchmarks, targeting distinct visual-linguistic reasoning abilities:
\begin{itemize}
    \item \textbf{General VQA Benchmarks}: % RealworldQA~\citep{grok}, MMBench$_\text{EN}$/MMBench$_\text{EN}$~\citep{liu2024mmbench}, SEED-Bench~\citep{li2023seed} — cover daily visual question answering with diverse scenes and question types, but performance is prone to fluctuations due to real-world variability.
    RealworldQA~\citep{grok} focuses on daily visual question answering with diverse scenes and question types; MMBench$_\text{EN}$~\citep{liu2024mmbench} is an English multimodal benchmark covering various visual-linguistic tasks; SEED-Bench~\citep{li2023seed} consists of 19K human-annotated multiple-choice questions spanning 12 evaluation dimensions including image understanding and generative comprehension. % Similar to other general VQA benchmarks, their performance is prone to fluctuations due to real-world variability.
    \item \textbf{Real-World Reasoning Benchmarks}: %AI2D~\citep{kembhavi2016diagram}, MMMU~\citep{yue2024mmmu}, MathVista~\citep{lu2023mathvista} — focus on complex multi-step reasoning (diagram interpretation, cross-disciplinary problem-solving), leading to large performance fluctuations across model variants.
    AI2D~\citep{kembhavi2016diagram} specializes in multi-step reasoning for diagram interpretation; MMMU~\citep{yue2024mmmu} is an interdisciplinary multimodal understanding and reasoning benchmark covering professional knowledge across domains; MathVista~\citep{lu2023mathvista} evaluates mathematical reasoning capabilities in visual contexts. % These benchmarks demand complex reasoning, leading to significant performance variations across model variants.
    \item \textbf{Text-Rich Benchmarks}: % ChartQA~\citep{liu2024chatqa}, TextVQA~\citep{singh2019towards}, OCRbench~\citep{Liu_2024} — specialize in text-related tasks (chart analysis, scene text understanding, OCR). These benchmarks exhibit the highest stability due to unambiguous task definitions and consistent annotations, serving as the primary basis for observing experimental phenomena and verifying core design effectiveness.
    ChartQA~\citep{liu2024chatqa} focuses on textual semantic understanding and question answering for chart data; TextVQA~\citep{singh2019towards} tests the ability to recognize and understand scene text; OCRbench~\citep{Liu_2024} comprehensively assesses model performance on OCR-related tasks. With unambiguous task definitions and consistent annotations, these benchmarks exhibit the highest stability and serve as the primary basis for verifying core design effectiveness.
\end{itemize}

\subsubsection{Image Generation Benchmarks}
We select three complementary benchmarks to comprehensively evaluate text-to-image alignment, knowledge grounding, and complex semantic modeling:
\begin{itemize}
    \item \textbf{GenEval}~\citep{ghosh2023geneval}: An object-focused framework assessing fine-grained text-image alignment (\textit{e.g.}, count, color, position, attributes) with high human annotation agreement, distinguishing simple/complex prompt performance.
    \item \textbf{WISE}~\citep{niu2025wise}: A world knowledge-informed benchmark verifying if generated images comply with common sense (\textit{e.g.}, cultural norms, physical laws, biology), focusing on knowledge grounding.
    \item \textbf{DPG-Bench}~\citep{hu2024ella}: Evaluates complex semantic alignment and instruction-following via dense structured prompts, testing entity recognition, attribute modeling, and relational reasoning with ``hard prompts'' for model differentiation.
\end{itemize}


\subsection{Image Understanding Performance} 
We evaluate the understanding performance of Kelix on diverse perspectives in Table \ref{tab:full_comprehensive_evaluation_clean}. Our Kelix (8B) showcases clear advantages when compared to other unified multimodal LLMs, regardless of whether they employ discrete, continuous, or hybrid tokenization schemes. 

Among discrete-tokenization unified models, Kelix achieves a decisive performance edge. For instance, it outperforms Janus-Pro (7B) by a substantial margin of 28.4 points on OCRBench (86.7 vs. 59.0) and surpasses Emu3 (8B) by 29.0 points on MathVista (76.5 vs. 47.6), demonstrating remarkable improvements on both text-rich and reasoning-intensive tasks.

Notably, Kelix narrows or even eliminates the long-standing performance gap with models built on continuous or hybrid tokenization, which are conventionally regarded as possessing superior visual representational fidelity. Specifically, Kelix attains a MathVista score of 76.5, surpassing the 14B Bagel (73.1) and the 30B Manzano (73.3) despite its considerably more compact parameter scale. Furthermore, in text-rich scenarios, Kelix attains state-of-the-art performance among all open-source unified multimodal models; its OCRBench score of 86.7 even matches or surpasses that of proprietary models such as GPT-4o (73.6) and Gemini-2.5-Pro (86.2). These robust results solidly establish Kelix as a powerful and versatile unified framework for general multimodal understanding.

\paragraph{Kelix overcomes the information bottleneck.}
Furthermore, existing discrete-tokenization unified models typically suffer substantial performance degradation on text-rich benchmarks, a critical limitation stemming from the inherent information bottleneck of single-token quantization approaches, as discussed in Sec. ~\ref{sec:discrete_image_tokenizer}. For instance, X-Omni (7B), the second-best discrete-tokenization baseline, achieves an OCRBench score of 70.4 and a TextVQA score of 77.4, which are notably lower than understanding-only Qwen2.5-VL (7B) (86.4 on OCRBench and 84.9 on TextVQA). % This gap highlights a prevalent challenge in existing discrete-token frameworks: single-token quantization fails to preserve sufficient textual and visual-semantic information for text-intensive tasks.

By contrast, Kelix stands as the first discrete-tokenization unified model to achieve OCRBench performace on par with that of understanding-only VLMs. Specifically, Kelix attains an OCRBench score of 86.7, which is nearly identical to the 86.4 achieved by VLM Qwen2.5-VL (7B), while outperforming the second-best discrete-tokenization baseline X-Omni (7B) by a substantial 16.3 points, corresponding to a relative performance improvement of 23\%. The result validates that Kelix’s multi-token quantization mechanism effectively bridges the representational gap between continuous visual embeddings and discrete tokens, demonstrating that discrete visual representations are no longer a fundamental performance bottleneck for multimodal understanding even for text-rich tasks that demand fine-grained semantic preservation.

% By contrast, Kelix represents the first discrete-tokenization unified model to match the text-rich benchmark performance of understanding-only VLMs. Specifically, OCRBench score of Kelix reaches 86.7, which is nearly identical to that of Qwen2.5VL (7B) (86.4), and outperforms the second best discrete-tokenization, i.e., X-Omni (7B) by 16.3, accomplishing a relative performance gain of 23\%. 
% The result validates that Kelix’s multi-token quantization mechanism effectively bridges the representational gap between continuous visual embeddings and discrete tokens, demonstrating that discrete visual representations are no longer a fundamental performance bottleneck for multimodal understanding, even for text-rich tasks that demand fine-grained semantic preservation.

% This result validates that Kelix’s multi-token quantization mechanism effectively bridges the representational gap between continuous visual embeddings and discrete tokens, demonstrating that discrete visual representations are no longer a fundamental performance bottleneck for multimodal understanding—even on text-heavy tasks that demand fine-grained semantic preservation.

% Futhermore, Kelix is the first model that achieves the OCRBench performace comparable to understanding-only models of the same model scale, e.g., Qwen2.5VL (86.7 v.s. 86.4). The result shows that Kelix can effectively close the representation ability between continuous embeddings and discrete tokens through multi-token quantization and demonstrates the discrete representation is no longer the bottleneck of multimodal understanding.

\subsection{Image Generation Performance} 
% Tables \ref{tab:generation_results} and \ref{tab:dpg_bench_results} summarize the image generation performance. On GenEval, Kelix (8B) achieves a leading score of 87.6, outperforming all discrete-tokenization models (e.g., Janus-Pro and Emu3) and even the 14B Bagel (82.0). It also leads the WISE benchmark with a score of 57.0, surpassing larger models like Manzano (30B, 54.0). Further results on DPG-Bench (85.5) confirm its robust compositional generation and text-alignment capabilities. 
% Tables \ref{tab:generation_results} and \ref{tab:dpg_bench_results} summarize the image generation performance of Kelix and comparative models. On the GenEval benchmark, Kelix (8B) achieves a state-of-the-art (SOTA) score of 87.6 among discrete-tokenization methods, which is significantly superior to the second-place models in this category, i.e., Nextflow (7B, 83.0) and X-Omni (7B, 83.0). Also, the GenEval score even outperforms large-scale 27B Qwen-Image by 0.6 (87.6 v.s. 87.0). Even when expanding the comparison to all methods including continuous tokenization, Kelix’s performance is only second to TUNA (7B, 90.0), demonstrating its strong competitiveness across paradigms. 

% For the WISE benchmark, Kelix attains a score of 57.0, ranking second only to Nextflow (7B, 59.0) while outperforming all continuous-tokenization models, (including larger model Bagel 14B, 52.0), and dedicated text-to-image (T2I) models (including larger model FLUX.1-dev 12B, 50.0). On DPG-Bench, Kelix achieves a score of 85.5, which is closely comparable to the SOTA performance of X-Omni (7B, 87.7) with a minimal gap, confirming its robust capabilities in compositional generation and fine-grained text-image alignment.  Moreover, our model does not have a reinforcement learning training phase like X-Omni does, demonstrating that our Kelix can well incorporate different modalities, striking a good balance between image understanding tasks and image generation tasks.

Tables \ref{tab:generation_results} and \ref{tab:dpg_bench_results} present the image generation performance of Kelix against competitive baselines across representative benchmarks. 

On the GenEval benchmark, Kelix (8B) attains a state-of-the-art (SOTA) score of 87.6 for discrete-tokenization unified models, a result that significantly outperforms the second-ranked models in this paradigm—Nextflow (7B, 83.0) and X-Omni (7B, 83.0). Notably, this score even exceeds that of the large-scale Qwen-Image (27B), which reaches 87.0, by a 0.6-point margin. When extending the comparison to all models across continuous, hybrid, and discrete tokenization paradigms, Kelix’s performance is second only to TUNA (7B, 90.0). Specifically, Kelix outperforms the second best discrete-tokenization unified models by substantial margins across several sub-tasks: 6 points in Counting, 6 points in Position, and 10 points in Color attributes. % further highlighting its strong cross-paradigm competitiveness in fine-grained text-image alignment.


For the WISE benchmark, which evaluates world knowledge grounding in generated images, Kelix attains a score of 57.0, ranking second only to Nextflow (7B, 59.0). Critically, Kelix outperforms all continuous-tokenization-based unified models, including the larger-scale Bagel (14B, 52.0), as well as dedicated text-to-image (T2I) models such as the larger FLUX.1-dev (12B, 50.0). On DPG-Bench, which measures compositional generation and complex semantic alignment capabilities, Kelix achieves a score of 85.5, which is highly competitive with the SOTA performance of X-Omni (7B, 87.7) with only a slight 2.2-point gap. The result confirms the robust capabilities of Kelix in image generation and fine-grained text-image alignment.
% For the WISE benchmark, which evaluates world knowledge grounding in generated images, Kelix attains a score of 57.0, ranking second only to Nextflow (7B, 59.0). Critically, Kelix outperforms all continuous-tokenization-based unified models—including the larger-scale Bagel (14B, 52.0)—as well as dedicated text-to-image (T2I) models such as the larger FLUX.1-dev (12B, 50.0). On DPG-Bench, which measures compositional generation and complex semantic alignment capabilities, Kelix achieves a score of 85.5, which is highly competitive with the SOTA performance of X-Omni (7B, 87.7) with only a slight 2.2-point gap, confirming its robust capabilities in compositional generation and fine-grained text-image alignment. 

Notably, Kelix is trained without the reinforcement learning stage adopted in some other baselines, e.g., X-Omni, which further validates that our framework can strike a good balance between multimodal understanding and high-quality image generation capabilities within a fully discrete autoregressive architecture.

% We present the quantitative results for our model’s image generation capabilities, evaluating them on two benchmarks: GenEval, WISE and DPG. As shown in Table~\ref{tab:generation_results} and Table~\ref{tab:dpg_bench_results}, while there remains a performance gap between our model and large-scale specialized models (e.g., FLUX.1) or proprietary systems like GPT-4o, our approach demonstrates competitive performance within the category of unified multimodal LLMs. Notably, on the WISE benchmark, which requires deep world knowledge, our model (43.0) outperforms established unified models such as Harmon-1.5B (41.0) and Janus-Pro-7B (35.0). Furthermore, our model maintains a balanced performance across both instruction following (GenEval) and semantic grounding (WISE), matching the overall WISE score of the much larger SDXL-3.5B. These results suggest that our unified architecture and training strategy provide a solid baseline for text-to-image generation, effectively capturing complex attributes and world knowledge even with a more compact parameter budget.



\section{Ablation Experiments}
\subsection{Ablation Study on Kelix-Tok}



\definecolor{colA}{RGB}{230, 242, 255} % Group A: 天空蓝 (Clear Sky Blue)
\definecolor{colB}{RGB}{230, 250, 245} % Group B: 薄荷绿 (Soft Mint)
\definecolor{colC}{RGB}{255, 248, 220} % Group C: 奶黄 (Cornsilk Yellow)
\definecolor{colD}{RGB}{255, 235, 225} % Group D: 杏粉 (Pale Apricot)
\definecolor{colE}{RGB}{255, 230, 240} % Group E: 玫瑰粉 (Soft Rose)
\definecolor{colF}{RGB}{240, 235, 255} % Group F: 薰衣草紫 (Light Lavender)
\definecolor{colG}{RGB}{240, 240, 240} % Group G: 中性灰 (Neutral Gray)
% ---------------------------------------------------

\begin{table*}[t]
    \centering
    \scriptsize 
    \setlength{\tabcolsep}{2.5pt} 
    \renewcommand{\arraystretch}{1.25} 
    
    \caption{
        Ablation studies on \textbf{Kelix-Tok}. 
        %``N/C'' denotes non-convergence.
    }
    \label{tab:ablation_MT_tokenizer}
    
    \begin{tabular}{
        @{} l  
        cccc       % General
        cccc       % Knowledge
        cccc       % Text-rich
        @{}
    }
        \toprule
        \multirow{2.5}{*}{\textbf{Configuration}} & 
        \multicolumn{4}{c}{\textbf{General Benchmarks}} & 
        \multicolumn{4}{c}{\textbf{Knowledge Benchmarks}} & 
        \multicolumn{4}{c}{\textbf{Text-rich Benchmarks}} \\
        
        \cmidrule(r){2-5} \cmidrule(lr){6-9} \cmidrule(l){10-13}
        
        & RWQA & MMB$_\text{CN}$ & MMB$_\text{EN}$ & \multicolumn{1}{c}{\textbf{Avg.}} 
        & AI2D & MMMU & MathV & \multicolumn{1}{c}{\textbf{Avg.}}
        & Chart & Text & OCRB & \multicolumn{1}{c}{\textbf{Avg.}} \\
        \midrule
        
        % ==========================================
        % Group (a)
        \multicolumn{13}{c}{\textit{(a) Baseline Comparison}} \\
        \cdashline{1-13}
        \rowcolor{colA}
        Discrete Baseline
        & 55.3 & 53.8 & 57.8 & 55.6
        & 61.6 & 31.3 & 40.3 & 44.4
        & 68.7 & 64.8 & 66.0 & 66.5 \\
        \rowcolor{colA}
        Continuous Baseline
        & 59.7 & 56.6 & 63.1 & 59.8
        & 65.4 & 34.2 & 50.3 & 50.0
        & 78.9 & 75.8 & 79.3 & 78.0 \\
        \rowcolor{colA}
        Kelix-Tok (Default)
        & 58.8 & 55.9 & 59.3 & 58.0
        & 62.9 & 33.9 & 50.0 & 48.9
        & 76.5 & 74.7 & 79.2 & 76.8 \\
        
        \specialrule{\lightrulewidth}{0pt}{0pt} 
        
        % % ==========================================
        % % Group (b)
        % \multicolumn{13}{c}{\textit{(b) Number of Sub-Tokens ($N$, Mean Pool)}} \\
        % \cdashline{1-13}
        % \rowcolor{colB}
        % $N=1$
        % & 55.3 & 53.8 & 57.8 & 55.6
        % & 61.6 & 31.3 & 40.3 & 44.4
        % & 68.7 & 64.8 & 66.0 & 66.5 \\
        % \rowcolor{colB}
        % $N=2$
        % & 57.6 & 56.2 & 57.7 & 57.2
        % & 62.7 & 39.5 & 46.4 & 49.5
        % & 74.2 & 67.9 & 71.7 & 71.3 \\
        % \rowcolor{colB}
        % $N=4$
        % & 61.7 & 57.8 & 62.2 & 60.6
        % & 65.8 & 39.9 & 46.1 & 50.6
        % & 75.4 & 72.9 & 75.0 & 74.4 \\
        % \rowcolor{colB}
        % $N=8$ (Default)
        % & 57.4 & 56.3 & 59.3 & 57.7
        % & 64.6 & 38.5 & 48.6 & 50.6
        % & 76.4 & 73.3 & 76.5 & 75.4 \\
        % \rowcolor{colB}
        % $N=16$ 
        % & 59.5 & 60.4 & 64.3 & 61.4
        % & 65.7 & 38.3 & 46.6 & 50.2
        % & 77.5 & 73.6 & 78.4 & 76.5 \\
        
        % \specialrule{\lightrulewidth}{0pt}{0pt} 

        % ==========================================
        % Group (c)
        \multicolumn{13}{c}{\textit{(b) Independent vs. Shared Codebooks (fixed $N=8$, Mean Pool)}} \\
        \cdashline{1-13}
        \rowcolor{colC}
        Shared Codebook
        & 49.0 & 46.8 & 50.6 & 48.8
        & 59.0 & 34.8 & 44.3 & 46.0
        & 67.0 & 60.4 & 61.7 & 63.0 \\
        \rowcolor{colC}
        Independent (Default)
        & 57.4 & 56.3 & 59.3 & 57.7
        & 64.6 & 32.4 & 48.3 & 48.4
        & 76.4 & 73.3 & 76.5 & 75.4 \\
        
        \specialrule{\lightrulewidth}{0pt}{0pt}

        % ==========================================
        % Group (d)
        \multicolumn{13}{c}{\textit{(c) Fusion Strategy (fixed $N=8$)}} \\
        \cdashline{1-13}
        \rowcolor{colD}
        Mean Pooling
        & 57.4 & 56.3 & 59.3 & 57.7
        & 64.6 & 32.4 & 48.3 & 48.4
        & 76.4 & 73.3 & 76.5 & 75.4 \\
        \rowcolor{colD}
        Sum Pooling (Default)
        & 58.8 & 55.9 & 59.3 & 58.0
        & 62.9 & 33.9 & 50.0 & 48.9
        & 76.5 & 74.7 & 78.5 & 76.6 \\
        
        \specialrule{\lightrulewidth}{0pt}{0pt}

        % ==========================================
        % Group (e)
        \multicolumn{13}{c}{\textit{(d) Codebook Dimension (d)}} \\
        \cdashline{1-13}
        \rowcolor{colE}
        $d=32$
        & 54.9 & 49.4 & 54.6 & 53.0
        & 60.1 & 30.7 & 43.3 & 44.7
        & 69.1 & 62.4 & 64.0 & 65.2 \\
        \rowcolor{colE}
        $d=64$
        & 53.3 & 52.6 & 56.6 & 54.2
        & 61.0 & 32.0 & 42.2 & 45.1
        & 69.4 & 65.4 & 65.8 & 66.9 \\
        \rowcolor{colE}
        $d=128$ (Default)
        %& 57.4 & 56.3 & 59.3 & 57.7
        %& 64.6 & 32.4 & 48.3 & 48.4
        %& 76.4 & 73.3 & 76.5 & 75.4 \\
        & 55.3 & 53.8 & 57.8 & 55.6
        & 61.6 & 31.3 & 40.3 & 44.4
        & 68.7 & 64.8 & 66.0 & 66.5 \\
        \rowcolor{colE}
        $d=256$
        & 52.0 & 50.9 & 51.9 & 51.6
        & 59.8 & 29.0 & 44.1 & 44.3
        & 71.1 & 63.3 & 63.8 & 66.1 \\
        \rowcolor{colE}
        $d=512$
        & 50.6 & 52.9 & 56.3 & 53.3
        & 59.9 & 29.6 & 44.8 & 44.8
        & 69.8 & 66.6 & 66.7 & 67.7 \\
        \rowcolor{colE}
        $d=1024$
        & 49.5 & 52.6 & 54.2 & 52.1
        & 60.1 & 32.0 & 41.7 & 44.6
        & 68.0 & 66.3 & 64.4 & 66.2 \\
        
        \specialrule{\lightrulewidth}{0pt}{0pt} 
        
        % Group (G)
        \multicolumn{13}{c}{\textit{(e) Codebook Size (\textit{S})}} \\
        \cdashline{1-13}
        \rowcolor{colF}
        \textit{S}=32,768
        & 54.0 & 50.1 & 53.9 & 52.7
        & 60.3 & 31.0 & 42.6 & 44.6
        & 70.0 & 64.0 & 65.0 & 66.3 \\
        \rowcolor{colF}
        \textit{S}=65,536 (Default)
        & 55.3 & 53.8 & 57.8 & 55.6
        & 61.6 & 31.3 & 40.3 & 44.4
        & 68.7 & 64.8 & 66.0 & 66.5 \\
        \rowcolor{colF}
        \textit{S}=131,072
        & 54.0 & 48.4 & 51.9 & 51.4
        & 59.6 & 32.6 & 43.8 & 45.3
        & 70.2 & 65.1 & 65.9 & 67.1 \\
        \bottomrule
    \end{tabular}
\end{table*}






% 下面强调了我们训的是stage1 & 3
In this section, we conduct a comprehensive set of ablation studies to systematically evaluate the impact of key design decisions in the Kelix-Tok. Table~\ref{tab:ablation_MT_tokenizer} summarizes our results. For efficiency, all model variants are trained on a representative subset of the full dataset and only adopt the two-stage training protocol (Stage 1 and Stage 3) — distinct from the main experiment’s complete training workflow. Each experiment follows a single-factor ablation protocol: starting from the Discrete Baseline, we modify one component at a time to precisely isolate its effect on performance.

\textbf{Baseline Comparison.} The results in Table~\ref{tab:ablation_MT_tokenizer}(a) reveals a critical trade-off among different visual representation schemes. We observe that, compared to the Continuous Baseline which represents a performance upper bound, the Discrete Baseline introduces a significant performance gap. This is particularly pronounced on Text-rich tasks, where its performance drops to approximately 85\% of the continuous baseline (66.5 vs. 78.0 Avg.) due to the information bottleneck. 
Our Kelix-Tok, leveraging a multi-token combinatorial representation, successfully achieves discretization while attaining performance comparable to its continuous counterpart.


% \textbf{Number of Sub-Tokens.}
% We examine the effect of varying the number of sub-tokens, $N$. As shown in Table~\ref{tab:ablation_MT_tokenizer}(b), increasing $N$ does not consistently yield significant improvements on understanding- and reasoning-intensive tasks (General and Knowledge). However, on Text-rich benchmarks—which require heightened sensitivity to visual detail—there is a clear positive correlation between $N$ and performance: the average score rises from 71.2 at $N=2$ to 76.5 at $N=16$. This pattern suggests that a larger number of sub-tokens primarily boosts the model's fine-grained perceptual capacity. To balance overall task performance with the computational cost introduced by additional tokens, we select $\mathbf{N=8}$ as the optimal compromise between effectiveness and efficiency.


\textbf{Independent vs. Shared Codebooks.}
We validate the necessity of our independent sub-codebook design. As shown in Table~\ref{tab:ablation_MT_tokenizer}(b), forcing all sub-tokens to a single Shared Codebook, as opposed to our default Independent codebooks, result in a catastrophic decline in performance (e.g., -12.3\% on average for Text-rich tasks). This finding highlights that sub-codebooks are essential for learning disentangled and meaningful visual concepts; utilizing a shared semantic space induces severe representation confusion.  These results underscore the critical role of our independent codebook strategy in achieving effective visual tokenization.

\textbf{Fusion Strategy.}
With $N=8$, we compare Mean Pooling and Sum Pooling as strategies for fusing sub-token embeddings. As shown in Table~\ref{tab:ablation_MT_tokenizer}(c), Sum Pooling consistently achieves equal or superior performance across all tasks, with a notable advantage on detail-sensitive Text-rich benchmarks (76.6 vs. 75.4 Avg.). This supports our hypothesis that summation more effectively preserves the overall information content and signal magnitude of the combinatorial representation than averaging, thus more effectively conveying visual details to the language model. Consequently, we adopt Sum Pooling as our default fusion strategy.

\textbf{Codebook Dimension.}
As shown in Table~\ref{tab:ablation_MT_tokenizer}(d), performance peaks at $d=128$. Smaller dimensions ($d \le 64$) lead to performance degradation due to an information bottleneck, while larger dimensions ($d \ge 256$) likely hinder optimization by introducing excessive sparsity in the representation space. Therefore, $d=128$ provides an optimal balance between representational capacity and optimization tractability, and is adopted as our default setting.

% \textbf{Initialization Strategy.}

\textbf{Codebook Size.}
Finally, we investigate the impact of codebook size on the discrete baseline. Results in Table~\ref{tab:ablation_MT_tokenizer}(e) show that increasing the size($S$) from 32K to 64K yields consistent gains. However, a further increase to 128K leads to performance degradation, likely due to insufficient training of an overly large codebook with limited data (i.e., codebook under-utilization). We thus select 65,536 as the optimal size, balancing representational capacity with training efficiency.

\textbf{Quantization Strategies.}
To investigate the impact of different discretization schemes, we compare \textbf{(1)} \textbf{Vector Quantization (VQ)} ~\citep{van2017neural,zhu2025addressing} with a codebook size of 65536, \textbf{(2)} \textbf{Finite Scalar Quantization (FSQ)}~\citep{mentzer2023finite} with an equivalent index space ($16^4 = 65536$), and \textbf{(3)} \textbf{Residual Quantization (RQ)}~\citep{lee2022autoregressive} utilizing a three-layer hierarchical structure ($3 \times 8192$). 
OCRBench is employed to evaluate their performance in fine-grained text recognition and document understanding. 
We observe that \textbf{VQ achieves the highest score of 66.0}, outperforming RQ (60.3) and FSQ (49.2). This performance gap indicates that the learnable embeddings in VQ offer superior representational flexibility for capturing dense visual-textual features. 
Consequently, VQ is adopted as the default quantization strategy for our final model.









% 导言区（需添加这两行配置子标题居中）
\captionsetup[subfigure]{justification=centering}  % 全局强制子图标题居中
\captionsetup[figure]{justification=centering}     % 新增：全局强制figure主标题居中

% 并排展示两张图片的figure环境
\begin{figure}[t!]
  \centering  % 整体居中
  % 第一张子图
  \begin{subfigure}{0.4775\textwidth}
    \centering  % 确保图片和子标题都居中
    \includegraphics[width=\textwidth]{OneBase/figures/stage3_loss_comparison.png}
    \subcaption{LM Loss curves on $N$ ablation}  % 自动生成(a)标签，无需手动加
    \label{fig:N-ab-loss}
  \end{subfigure}
  \hspace{0.5mm}  % 极致小的固定间距（0.5毫米）
  % 第二张子图
  \begin{subfigure}{0.475\textwidth}
    \centering  % 确保图片和子标题都居中
    % \includegraphics[width=\textwidth]{Helix/figures/stage3_textrich_average_curve.png}
    % \subcaption{Text-rich bench average curves on $N$ ablation}  % 自动生成(b)标签，无需手动加
    \includegraphics[width=\textwidth]{OneBase/figures/stage3_ocr_bench_score_curve.png}
    \subcaption{OCRBench curves on $N$ ablation}
    \label{fig:N-ab-ocr}
  \end{subfigure}
  
  \caption{Ablation study on parameter $N$: LM Loss and OCRBench curves.}
  \label{fig:N-ab}
\end{figure}

\textbf{Number of Sub-Tokens.} We investigate the impact of varying the number of sub-tokens $N$ by conducting experiments with $N=1,2,4,8,16$ alongside a continuous baseline, with results presented in Fig.~\ref{fig:N-ab}. As shown in Subfigure.~\ref{fig:N-ab-loss}, the LM loss consistently decreases with increasing $N$, and the loss curves for $N \geq 8$ converge closely to the continuous baseline, indicating that sufficient discrete tokens can effectively compress continuous visual information without substantial loss. We further evaluate our training checkpoints on the OCRBench benchmark, which are highly sensitive to model performance.  Subfigure.~\ref{fig:N-ab-ocr} demonstrates that OCRBench scores consistently increase with $N$. Upon training convergence, the performance of the continuous baseline is nearly matched by the top two discrete configurations ($N=8$ and $N=16$), which further validates the robustness of our multi-token based Kelix.

% As shown in Table~\ref{tab:ablation_MT_tokenizer}(b), increasing $N$ does not consistently yield significant improvements on understanding- and reasoning-intensive tasks (General and Knowledge). However, on Text-rich benchmarks—which require heightened sensitivity to visual detail—there is a clear positive correlation between $N$ and performance: the average score rises from 71.2 at $N=2$ to 76.5 at $N=16$. This pattern suggests that a larger number of sub-tokens primarily boosts the model's fine-grained perceptual capacity. To balance overall task performance with the computational cost introduced by additional tokens, we select $\mathbf{N=8}$ as the optimal compromise between effectiveness and efficiency.




\subsection{Ablation Study on Kelix-LLM}
\begin{table*}[t]
    \centering
    \scriptsize 
    \setlength{\tabcolsep}{2.5pt} 
    \renewcommand{\arraystretch}{1.25} 
    
    \caption{
        Ablation studies on our \textbf{Unified Multimodal LLM Architecture}. 
        We evaluate the design of multimodal interfaces, representation paradigms, and training synergies across Stage 1 and Stage 3. 
    }
    \label{tab:ablation_unified_model}
    
    \begin{tabular}{
        @{} l  
        cccc       % General
        cccc       % Knowledge
        cccc       % Text-rich
        @{}
    }
        \toprule
        \multirow{2.5}{*}{\textbf{Configuration}} & 
        \multicolumn{4}{c}{\textbf{General Benchmarks}} & 
        \multicolumn{4}{c}{\textbf{Knowledge Benchmarks}} & 
        \multicolumn{4}{c}{\textbf{Text-rich Benchmarks}} \\
        
        \cmidrule(r){2-5} \cmidrule(lr){6-9} \cmidrule(l){10-13}
        
        & RWQA & MMB$_\text{CN}$ & MMB$_\text{EN}$ & \multicolumn{1}{c}{\textbf{Avg.}} 
        & AI2D & MMMU & MathV & \multicolumn{1}{c}{\textbf{Avg.}}
        & Chart & Text & OCRB & \multicolumn{1}{c}{\textbf{Avg.}} \\
        \midrule
        
        % ==========================================
        % Group (a) Multimodal Interface Design
        % \multicolumn{13}{c}{\textit{(a) Design of Multimodal Input/Output Interface}} \\
        %\cdashline{1-13}
        %\rowcolor{colA}
        exp1: Unified Baseline
        & 61.3 & 67.7 & 66.9 & 65.3 &
        70.8	& 42.0 &	55.8 &	56.2 &
        76.9 &	72.3	& 79.9 &	76.4 \\
        \rowcolor{colA}
        exp2: Random Init Input Emb.
        & 56.7 &	60.7&	59.3&	58.9	
        &67.5&	34.7&	53.8&	52.0	
        &77.8	&70.4&	77.0 &	75.1 \\
        \rowcolor{colA}
        exp3: 9-Head Predictor (Output)
        & 64.8	&60.4	&61.6	&62.3 
        &	69.3	&41.7&	52.9 &	54.6 
        &	75.2 &	71.4 &	76.5 &	74.4  \\

        \bottomrule
    \end{tabular}
\end{table*}

In this section, we ablate the key design components in our unified LLM architecture, including the multimodal input/output interface, representation paradigm, and training data synergy, to validate their effectiveness. The experimental setting are aligned with the settings in Table~\ref{tab:ablation_unified_model}, where the exp1 served as the baseline and modify one core component at a time to isolate its impact on model performance. In the following, we detail the training setting and analyze the the experimental results, and derive insights for the unified LLM framework.

\paragraph{Experimental Setting Explanations.}


All ablation variants are built on the Qwen3-8B backbone but only adopt a two-stage training protocol (Stage 1 and Stage 3) with a subset of full training data, unlike the full four-stage workflow of the main experiment. This simplification streamlines the ablation process by eliminating redundant training stages that are unnecessary for isolating the effects of target components. It ensures the ablation focuses on core design validations without excessive computational overhead. Key ablations are defined as follows:
\begin{itemize}

\item \textbf{exp1: Unified Baseline}: The full-version of our proposed unified LLM model. It includes: (1) Stage 1 training with a Linear Aligner to map VQ discrete tokens to the LLM's semantic space; (2) a unified \textit{Block Encoder} (sum pooling) and \textit{Block Decoder} for multimodal input aggregation and output reconstruction; (3) training on both understanding-oriented datasets (e.g., VQA, OCR) and generation-oriented datasets (e.g., text-to-image pairs). This model variant serves as the performance benchmark for all ablations.

\item \textbf{exp2: Random Init Input Emb.}: Ablates the \textit{up-projector} component in LLM Stage 1. Starting from LLM Stage 1, we directly use token IDs from the VQ quantizer as visual inputs, and initialize the visual embedding table randomly without any alignment training. This setting could test whether the Linear Aligner and Stage 1's cross-modal alignment strategy are necessary.

\item \textbf{exp3: 9-Head Predictor (Output)}: Replaces the \textit{Block Decoder} with a 9-head parallel predictor. Specifically, the LLM's single hidden state is projected into 9 independent vectors via 9 separate linear heads, which are then fed into the LM head to generate 9 discrete tokens. For text generation, only the first token is retained; for image generation, the latter 8 tokens are used. This variant verifies the rationality of the \textit{Block Decoder}'s hierarchical reconstruction mechanism.

\end{itemize}

\paragraph{Result Analysis}

We analyze the ablation results across three benchmark categories (General Benchmarks, Knowledge Benchmarks, Text-rich Benchmarks) and focus on key performance differences between variants.

The results in Table~\ref{tab:ablation_unified_model} highlight the critical role of the \textit{up-projector} and \textit{Block Encoder} in cross-modal modeling:

\begin{itemize}

\item \textbf{Effect of Linear Aligner (exp1 vs. exp2)}: Compared to the baseline (exp1, General Avg. 65.3), the variant using randomly initialized visual embedding table (exp2) shows a significant performance drop across all tasks: General Avg. decreases by 6.4 points (to 58.9), Knowledge Avg. drops by 4.2 points (56.2 → 52.0), and Text-rich Avg. decrease by 1.3 points (76.4 → 75.1). This confirms that Stage 1's \textit{up-projector} effectively bridges the semantic gap between discrete visual tokens and the LLM's text space, avoiding catastrophic forgetting of linguistic priors and laying the foundation for cross-modal reasoning. Without this alignment, random initialization of visual embeddings leads to misalignment between modalities, severely degrading model performance.

\item \textbf{Effect of \textit{Block Encoder} (exp1 vs. exp3) }: Replacing the \textit{Block Encoder} with a 9-head predictor (exp3) results in consistent performance degradation: General Avg. decreases by 3.0 points (65.3 → 62.3), Knowledge Avg. drops by 1.6 points (56.2 → 54.6), and Text-rich Avg. falls by 2.0 points (76.4 → 74.4). The 9-head design treats text and image outputs as independent token streams, breaking the semantic integrity of "blocks" (text blocks as single semantic units, image blocks as 8-token spatial clusters). This indicates that the \textit{Block Encoder}'s hierarchical reconstruction (aggregating hidden states into blocks and then reconstructing tokens) better preserves the structural and semantic consistency of multimodal data. In contrast, the parallel 9-head predictor introduces semantic fragmentation, especially for image tokens that require spatial coherence, thus reducing model performance.



\begin{figure}[t!]
  \centering
  \includegraphics[width=1\textwidth]{OneBase/figures/case_0210_v2.pdf}
    \caption{The visualization comparison of different unified models in image generation.}
    \label{fig:casecompare}
\end{figure}

\end{itemize}




\paragraph{Key Findings}

From the ablation study on the unified LLM architecture, we draw two critical conclusions:

\begin{enumerate}

\item \textbf{Cross-modal alignment is indispensable}: The \textit{up-projector} in Stage 1 effectively maps discrete visual tokens to the LLM's semantic space, avoiding inter-modal misalignment and ensuring the model's basic cross-modal reasoning ability.

\item \textbf{Unified block-wise processing is superior}: The \textit{Block Decoder}'s hierarchical reconstruction mechanism preserves the semantic integrity of multimodal data, outperforming parallel multi-head predictors that cause semantic fragmentation.

\end{enumerate}

These findings confirm the rationality of our unified LLM architecture design, providing empirical support for the core components (the up-projector and Block Embedding Moduling) and laying the foundation for the model's superior performance in both understanding and generation tasks.









\subsection{Case Study of Image Generation}

We show image generation cases of Kelix and other SOTA unified models in Figure~\ref{fig:casecompare}. The visualization results demonstrate that Kelix achieves superior instruction-following capability compared to baselines.
%
In the first case (pillowcase with ``DREAM BIG'' in cursive), the Janus-Pro, SD-mid, and TA-TOK produces either unclear or non-cursive text.
%
In the second case (garden with 3 red roses and 1 daisy), they generate incorrect flower counts or mix up the flower types.
%
In the third case (one black lamp, three white plates, and two red jars, a four-legged table), all of these baseline methods do not incorrectly render the number of plates or jars.
% Kolors
These comparisons highlight that Kelix better adheres to fine-grained instructions including text rendering, object counting, and spatial arrangement.






\section{Related Works}

We review prior work along the two representation dimensions introduced in Section~\ref{sec2_Design}: \emph{representation form} (discrete vs.\ continuous) and \emph{semantic level} (high-level vs.\ low-level), covering vision-language models, unified multimodal models, and image generation techniques.

\subsection{Vision-Language Models}

CLIP~\citep{radford2021learning} established the foundation for modern VLMs by aligning visual and textual semantics in a shared embedding space via large-scale contrastive pre-training. Building on CLIP's pre-trained ViT, Flamingo~\citep{alayrac2022flamingo} introduced cross-attention mechanisms to fuse visual and language representations, enabling few-shot visual reasoning. However, the cross-attention design modifies the LLM architecture, limiting compatibility with pre-trained LLM backbones.

BLIP-2~\citep{li2023blip} addressed this by proposing a lightweight Q-Former that bridges a frozen ViT and a frozen LLM, requiring only the adaptor to be trained. LLaVA~\citep{liu2023visual} further simplified the pipeline to a ``ViT--Projector--LLM'' framework, feeding fine-grained patch embeddings directly to the LLM. This paradigm has since become the dominant VLM design, adopted by Qwen2-VL, Qwen2.5-VL~\citep{wang2024qwen2, bai2025qwen2}, MiniCPM-V~\citep{yao2024minicpm}, InternVL~\citep{chen2024internvl}, among others.

More recently, VLMs have shifted toward native-resolution ViT encoders to process images at their original resolution without complex splitting or resizing, enabling more fine-grained visual understanding. Representative models include Qwen3-VL~\citep{Qwen3-VL}, Kimi-VL~\citep{team2025kimi}, Seed-VL -1.5~\citep{guo2025seed1}, and Keye-VL~\citep{yang2025kwai}. Note that all these VLMs rely on \emph{continuous} visual features and are limited to understanding tasks.

\subsection{Unified Multimodal Models}

Although understanding-only VLMs have achieved great success, the absence of generation capability limits both the application scope and the scaling potential of these models. This has motivated the community to develop unified models that support both understanding and generation. Existing approaches broadly fall into two categories: \emph{unified architectures} and \emph{hybrid architectures}.

\textbf{Unified Architectures.}
In unified architectures, a single representation serves both understanding and generation, and discrete tokenization is a natural choice for this paradigm. Early works such as Chameleon~\citep{team2024chameleon}, LlamaGen~\citep{sun2024autoregressive}, and Emu3~\citep{wang2024emu3} adopt VQ-VAE-like techniques~\citep{van2017neural} to encode images into low-level discrete codes, achieving unified text-vision modeling under next-token prediction. However, the semantic level mismatch between low-level visual codes and high-level text tokens leads to low cross-modal alignment efficiency and poor sample efficiency during scaling.

More recently, a growing number of works have turned to high-level visual semantics derived from ViT encoders, including BLIP3o~\citep{chen2025blip3ofamilyfullyopen}, BLIP3o-NeXT~\citep{chen2025blip3o}, TaTok~\citep{han2025vision}, X-Omni~\citep{geng2025xomnireinforcementlearningmakes}, UniTok~\citep{ma2025unitok}, and NextFlow~\citep{zhang2026nextflow}. High-level semantics exhibit significantly higher alignment efficiency, enabling fast transfer of the LLM's world knowledge and reasoning capabilities to visual generation, and these models have achieved strong results on generation benchmarks. However, discrete tokens in this route still face a fundamental challenge: the information loss inherent in discretization causes them to perform significantly worse than contemporary understanding-only VLMs on multimodal understanding tasks, which to some extent hinders the realization of true modal unification under this framework.

\textbf{Hybrid Architectures.}
To address the understanding weakness of unified architectures, hybrid models adopt separate representations for understanding and generation. On the understanding side, they retain the proven ViT--Projector--LLM pipeline from understanding-only VLMs; on the generation side, they employ a different representation for visual synthesis. Early hybrid models use low-level discrete codes for generation, such as Janus~\citep{wu2025janus}, Janus-Pro~\citep{janusPro} (VQ-VAE-based) and UniHetero~\citep{chen2025unihetero}, or low-level continuous features, such as JanusFlow~\citep{ma2025janusflow} and BAGEL~\citep{deng2025emerging}. Following the same trend as unified architectures, recent hybrid works have also begun exploring high-level semantics for the generation branch, including MetaQuery~\citep{pan2025transfer}, and Pisces~\citep{xu2025pisces}

The fundamental limitation of hybrid architectures is the representational asymmetry between understanding and generation, which inherently impairs knowledge sharing and cross-task transfer. Recent efforts such as Manzano~\citep{li2025manzano}, TuNa~\citep{liu2025tuna}, Ming-UniVision~\citep{mingUniVision} and RAE~\citep{tong2026scaling} have attempted to address this by unifying the representation space into high-level semantics, achieving some positive results.

\textbf{Our Position.}
Our proposed \myorange{\textbf{Kelix}} follows the \textbf{unified architecture} route with \textbf{high-level discrete semantics}, but introduces a novel \myorange{multi-token quantization mechanism} that expands the representational capacity of discrete tokens without increasing the LLM's sequence length. This directly addresses the information loss bottleneck of discrete tokenization, aiming to close the understanding gap between discrete and continuous representations and unlock the full potential of unified multimodal modeling.

\subsection{Image Generation}

Image Generation is undergoing a profound transformation. The core architecture of diffusion models is shifting from convolutional U-Net backbones~\citep{rombach2022high} toward more scalable and parameter-efficient Diffusion Transformers (DiT)~\citep{peebles2023scalable, team2025zimage}, with training objectives evolving from traditional noise prediction to more stable and expressive flow-matching paradigms~\citep{ma2025janusflow, xie2025sana}. A pivotal driver behind this evolution is the upgrade of conditioning mechanisms: transitioning from static, fixed-vocabulary text encoders like CLIP~\citep{radford2021learning} or T5 toward to more dynamic LLM-driven conditioning. As demonstrated by Qwen-Image~\citep{wang2024qwen2}, MetaQuery~\citep{pan2025transfer}, and Manzano~\citep{li2025manzano}, this paradigm harnesses the rich hidden states and world knowledge of pre-trained LLMs to guide the generative process. By treating the LLM as a sophisticated semantic engine, these systems transcend basic text-to-image mapping to unlock reasoning-heavy capabilities, such as \textbf{instruction-based editing}~\citep{qian2025pico} with nuanced spatial control and \textbf{in-context image generation}~\citep{sun2024generative} that maintains visual consistency. In this framework, image generation is effectively reimagined as a seamless extension of the LLM’s multimodal reasoning chain.


Our \textbf{Kelix} aligns with this trend but distinguishes itself by performing \textbf{semantic-level autoregressive modeling over high-level discrete tokens}. Unlike paradigms that merely utilize continuous hidden states as conditioning signals, Kelix adopts a modular \textit{Tokenizer--LLM--Detokenizer} design that fundamentally integrates generation into the LLM's next-token prediction framework~\citep{llamagen, wang2024emu3}. By predicting discrete visual representations~\citep{van2017neural}, Kelix allows the model to fully inherit the LLM’s logical capacity for complex generative tasks while leveraging a diffusion-based DiT as the \textbf{Detokenizer} to bypass the inherent quality limitations of traditional discrete decoders.

\section{Conclusion}
We present \textbf{Kelix}, a fully discrete autoregressive unified model that closes the understanding gap between discrete and continuous visual representations in multimodal models.
%
Our key insight is that a single discrete token cannot match the expressive capacity of a continuous embedding within a fixed vocabulary size. By decomposing each continuous patch embedding into multiple parallel discrete tokens, Kelix substantially expands the information capacity of discrete codes, enabling understanding performance comparable to continuous-feature VLMs.
%
To maintain efficiency despite the increased token count, we adopt a Next-Block Prediction paradigm that models grouped token blocks autoregressively, keeping the effective sequence length, and thus training and inference cost, on par with single-token approaches.
%
Following a modular \emph{Tokenizer--LLM--Detokenizer} design, each component can be developed and iterated independently: the tokenizer handles semantic alignment, the LLM performs reasoning and semantic control, and the diffusion-based detokenizer renders high-quality images.
%
Extensive experiments on major multimodal understanding and image generation benchmarks demonstrate state-of-the-art results among comparable-scale unified models, validating the effectiveness of our approach.
%
In the future, we plan to extend Kelix to video understanding and generation, moving toward a more general unified multimodal intelligence.


\bibliography{bibtex}
% \bibliographystyle{acm}
% \bibliographystyle{plainnat}
\bibliographystyle{IEEEtranN}


\appendix







\begin{figure}[t!]
  \centering
  \includegraphics[width=0.95\textwidth]{OneBase/figures/pdf_plotted_with_prompt_g1.jpg}
    \caption{The visualization of Kelix in image generation.}
    \label{fig:showcase}
\end{figure}

\begin{figure}[t!]
  \centering
  \includegraphics[width=0.95\textwidth]{OneBase/figures/pdf_plotted_with_prompt_g2.jpg}
    \caption{The visualization of Kelix in image generation.}
    \label{fig:showcase}
\end{figure}

\begin{figure}[t!]
  \centering
  \includegraphics[width=0.95\textwidth]{OneBase/figures/pdf_plotted_with_prompt_g3.jpg}
    \caption{The visualization of Kelix in image generation.}
    \label{fig:showcase}
\end{figure}





\newpage
\quad \\
\quad \\



\section{Contributions}
% Names marked with * denote individuals who have departed from our team. 

% \definecolor{damaiblue}{RGB}{0, 0, 100}
% \definecolor{damaigreen}{RGB}{0, 100, 0}
% \definecolor{damaired}{RGB}{100, 0, 0}
% \begin{multicols}{3}
% \noindent
% \textbf{ Core Contributors} \\
% Boyang Ding \\
% Chenglong Chu \\
% Dunju Zang \\
% Han Li \\
% Muhao Wei \\
% Jiangxia Cao \\
% Kun Gai \\
% Ruiming Tang \\
% Shiyao Wang \\
% Siyang Mao \\
% Xinchen Luo \\
% Yahui Liu \\
% Zhixin Ling \\
% Zhuoran Yang \\
% Ziming Li \\

% \noindent
% \textbf{ Contributors} \\
% Chengru Song \\ 
% Guowang Zhang \\
% Hao Peng \\
% Hao Wang \\
% Jin Ouyang \\
% Jinghao Zhang \\
% Lejian Ren \\
% Qiang Wang
% Tao Wang \\
% Xinmei Wang \\
% Ziqi Wang \\
% \end{multicols}


\noindent
\textbf{Core Contributors}\textsuperscript{\dag}\quad
Boyang Ding\textsuperscript{*},
Chenglong Chu,
Dunju Zang,
Han Li,
Jiangxia Cao,
Kun Gai,
Muhao Wei,
Ruiming Tang,
Shiyao Wang,
Siyang Mao,
Xinchen Luo,
Yahui Liu,
Zhixin Ling,
Zhuoran Yang,
Ziming Li

\vspace{0.5em}
\noindent
\textbf{Contributors}\textsuperscript{\dag}\quad
Chengru Song,
Guorui Zhou,
Guowang Zhang,
Hao Peng,
Hao Wang,
Jiaxin Deng,
Jin Ouyang,
Jinghao Zhang,
Lejian Ren,
Qianqian Wang,
Qigen Hu,
Tao Wang,
Xingmei Wang,
Yiping Yang,
Zixing Zhang,
Ziqi Wang

\vspace{0.5em}
\noindent
{\small \textsuperscript{\dag}Listed alphabetically by first name.\quad \textsuperscript{*}Project Leader.}


% \noindent
% \textbf{Core Contributors} \\[0.3em]
% \begin{tabular}{@{}llll@{}}
% Boyang Ding & Chenglong Chu & Dunju Zang & Han Li \\
% Jiangxia Cao & Kun Gai & Muhao Wei & Ruiming Tang \\
% Shiyao Wang & Siyang Mao & Xinchen Luo & Yahui Liu \\
% Zhixin Ling & Zhuoran Yang & Ziming Li & \\
% \end{tabular}

% \vspace{0.5em}
% \noindent
% \textbf{Contributors} \\[0.3em]
% \begin{tabular}{@{}llll@{}}
% Chengru Song & Guowang Zhang & Hao Peng & Hao Wang \\
% Jin Ouyang & Jinghao Zhang & Lejian Ren & Qiang Wang \\
% Tao Wang & Xinmei Wang & Ziqi Wang & \\
% \end{tabular}


\end{document}
